% Options for packages loaded elsewhere
\PassOptionsToPackage{unicode}{hyperref}
\PassOptionsToPackage{hyphens}{url}
\documentclass[
  ukrainian,
]{article}
\usepackage{xcolor}
\usepackage[margin=2.5cm]{geometry}
\usepackage{amsmath,amssymb}
\setcounter{secnumdepth}{-\maxdimen} % remove section numbering
\usepackage{iftex}
\ifPDFTeX
  \usepackage[T1]{fontenc}
  \usepackage[utf8]{inputenc}
  \usepackage{textcomp} % provide euro and other symbols
\else % if luatex or xetex
  \usepackage{unicode-math} % this also loads fontspec
  \defaultfontfeatures{Scale=MatchLowercase}
  \defaultfontfeatures[\rmfamily]{Ligatures=TeX,Scale=1}
\fi
\usepackage{lmodern}
\ifPDFTeX\else
\setmainfont{DejaVu Serif}
\setsansfont{DejaVu Sans}
\setmonofont{DejaVu Sans Mono}
\setmathfont{STIX Two Math}
\fi
\ifPDFTeX\else
  % xetex/luatex font selection
\fi
% Use upquote if available, for straight quotes in verbatim environments
\IfFileExists{upquote.sty}{\usepackage{upquote}}{}
\IfFileExists{microtype.sty}{% use microtype if available
  \usepackage[]{microtype}
  \UseMicrotypeSet[protrusion]{basicmath} % disable protrusion for tt fonts
}{}
\makeatletter
\@ifundefined{KOMAClassName}{% if non-KOMA class
  \IfFileExists{parskip.sty}{%
    \usepackage{parskip}
  }{% else
    \setlength{\parindent}{0pt}
    \setlength{\parskip}{6pt plus 2pt minus 1pt}}
}{% if KOMA class
  \KOMAoptions{parskip=half}}
\makeatother
\usepackage{graphicx}
\graphicspath{{/Users/kalinicenkovitalijmikolajovic/.codex/worktrees/5c15/Script_for_translate_book/Output_Final/}}
\setkeys{Gin}{width=0.82\linewidth,height=0.82\textheight,keepaspectratio}
\makeatletter
\newsavebox\pandoc@box
\newcommand*\pandocbounded[1]{% scales image to fit in text height/width
  \sbox\pandoc@box{#1}%
  \Gscale@div\@tempa{\textheight}{\dimexpr\ht\pandoc@box+\dp\pandoc@box\relax}%
  \Gscale@div\@tempb{\linewidth}{\wd\pandoc@box}%
  \ifdim\@tempb\p@<\@tempa\p@\let\@tempa\@tempb\fi% select the smaller of both
  \ifdim\@tempa\p@<\p@\scalebox{\@tempa}{\usebox\pandoc@box}%
  \else\usebox{\pandoc@box}%
  \fi%
}
% Set default figure placement to htbp
\def\fps@figure{htbp}
\makeatother
\ifLuaTeX
\usepackage[bidi=basic,shorthands=off,provide=*]{babel}
\else
\usepackage[bidi=default,shorthands=off,provide=*]{babel}
\fi
\ifLuaTeX
  \usepackage{selnolig} % disable illegal ligatures
\fi
\setlength{\emergencystretch}{3em} % prevent overfull lines
\providecommand{\tightlist}{%
  \setlength{\itemsep}{0pt}\setlength{\parskip}{0pt}}
\usepackage{bookmark}
\IfFileExists{xurl.sty}{\usepackage{xurl}}{} % add URL line breaks if available
\urlstyle{same}
\hypersetup{
  pdflang={uk},
  hidelinks,
  pdfcreator={LaTeX via pandoc}}

\author{}
\date{}

\begin{document}

\pandocbounded{\includegraphics[keepaspectratio]{images/_page_101_Picture_7.jpeg}}

Рис. 13.2. Схематичне представлення проблеми теплопровідності: a ---
тепло тече всередину, якщо ви(0, t) \textless{} 0, і назовні, якщо ви(0,
t) \textgreater{} 0; b --- резервуар настільки глибокий, що гранична
умова внизу не впливає на розв'язок при значеннях x.кістка дорівнює нулю
(орієнтир температур) (див. рис. 13.2). Наше завдання --- знайти
температурне поле в рідині на різних глибинах і в різний час. Іншими
словами, нам потрібно розв'язати проблему(13.6)

\begin{verbatim}
\begin{aligned}
(\text{УЧП})\quad & u_t = u_{xx}, && 0 < x < \infty,\; 0 < t < \infty, \\
(\Gamma \text{У})\quad & u_x(0,t) - u(0,t) = 0, && 0 < t < \infty, \\
(\text{НУ})\quad & u(x,0) = u_0, && 0 < x < \infty.
\end{aligned}
\end{verbatim}

Щоб розв'язати цю задачу, ми застосуємо перетворення Лапласа до змінної
t. Зверніть увагу, що можна також застосувати перетворення Лапласа до
змінної x, оскільки воно змінюється від 0 до \(\infty\) . Після
застосування перетворенняМи отримуємо звичайне диференціальне рівняння
за змінною X.(13.7)

\[
(ODY) \quad sU(x) - u_0 = \frac{d^2U}{dx^2}, \quad 0 < x < \infty,
(\Gamma Y) \quad \frac{dU}{dx}(0) = U(0)
\]

\begin{verbatim}
(Ми перетворили лише рівняння та граничну умову, а не початкову умову.) Маємо звичайне диференціальне рівняння другого порядку, але з однією граничною умовою при x=0. Насправді існує ще одна гранична умова. Він визначається фізичними міркуваннями і фіксується як U(x) \to 0 , коли x \to +\infty . Зверніть увагу, що для спрощення позначення ми опустили параметр s із аргументів функції U(x) всюди, оскільки рівняння (13.7) є диференціальним рівнянням лише відносно x.Щоб розв'язати задачу (13.7), запишіть загальне рішення (загальний розв'язок однорідного + часткового розв'язку гетерогенного) TAC:
\end{verbatim}

\[
U(x) = c_1 e^{V_{5x}} + c_2 e^{-V_{5x}} + u_0/s.
\]

Підставляючи цей вираз у граничну умову (13.7), можна визначити
константи \(c_1\) та \(c_2\) (одразу зрозуміло, що \(c_3 = 0\),

\pandocbounded{\includegraphics[keepaspectratio]{images/_page_102_Figure_7.jpeg}}

\begin{verbatim}
Рис. 13.3. Граф ймовірнісного інтеграла (eri) та його комплементарної функції (eric).Інакше температура буде зростати нескінченно зі збільшенням координати X). Визначаючи сталу c_2 з граничної умови при x=0, отримуємо фінальний вираз для U\left(x\right) (13.8)
\end{verbatim}

\[
U(x) = u_0 \left\{ \frac{e^{-V \widehat{s} x}}{s(V \widehat{s} + 1)} \right\} + \frac{1}{s}.
\]

\begin{verbatim}
Залишається зробити останній крок. Щоб визначити температурне поле, потрібно обчислити u\left(x,\,t\right)
\end{verbatim}

\[
u(x, t) = \mathcal{L}^{-1}[U(x, s)]
\]

(ми повернулися до написання всіх аргументів функції \(U(x, s)\)). Щоб
повернути перетворення Лапласа, ми використаємо таблиці, наведені в
додатку. У результаті ми отримуємо(13.9)

\[
u(x, t) = u_0 - u_0 \left[ \text{eric} \left( \frac{x}{2} \sqrt{t} \right) + \text{eric} \left( \sqrt{t} + \frac{x}{2} \sqrt{t} \right) e^{x+t} \right],
\]

де

\[
\operatorname{eric}(x) = \frac{2}{\sqrt{\pi}} \int_{x}^{\infty} e^{-\xi^{2}} d\xi
\]

є функцією, доповнюючою до ймовірнісного інтеграла. Графіки цієї функції
показані на рис. 13.3. Аналіз цього рішення P H C

\pandocbounded{\includegraphics[keepaspectratio]{images/_page_103_Figure_7.jpeg}}

13.4. Температура всередині напівнескінченного середовища в різний
час.Потребує дуже малих зусиль. Фіг. (13.4) представлено результати
чисельного аналізу, отриманого на комп'ютері, оснащеному плоттером.\#
ПРИМІТКИ- 1. Перетворення Лапласа також можна звести до неоднорідного
диференціального рівня з частковими множниками (для виконання методу
диференціації різних рівнів можна звести до одного рівня), але можна
звести його до рівня різних множників (метод диференціації можна звести
до рівня коефіцієнтів заміщення). Таблиця 1. 13.2 У сфері захисту шкіри
існує два методи.

\begin{itemize}
\tightlist
\item
  \begin{enumerate}
  \def\labelenumi{\arabic{enumi}.}
  \setcounter{enumi}{1}
  \tightlist
  \item
    Для виконання цих і межових завдань також проводиться трансформація
    Ганкеля і Мелліна, а також трансформація Лапласа в одному аспекті.
    Відтворення Лапласа підстановки оперативного множення формули
  \end{enumerate}
\end{itemize}

\[
\mathscr{L}[y'] = s\mathscr{L}[y] - y(0),
\]

Таблиця 13.2 Порівняння методів розділення перетворення Лапласа та
змінних\textbar{} \textbar{} Метод \textbar{} \textbar{}

\textbar-------------------------\textbar---------------------------\textbar----------------------------------\textbar{}

\textbar{} Перетворення Лапласа \textbar{} Розділення \textbar{}

Гетерогенні PPP \textbar{} Так \textbar{} Так \textbar{}

Гетерогенні ГІ \textbar{} Так \textbar{} Так \textbar{}

Змінні ймовірності \textbar{} Так \textbar{} de' \textbar{}

Нелінійні рівняння \textbar{} Так \textbar{} Так \textbar а перетворення
Ганкеля і Мелліна замінюють диференціальний оператор на множник.
Наприклад, для перетворення Ганкеля, визначеного формулою

\[
H[y] = \int_{0}^{\infty} r J_{o}(\xi r) y(r) dr,
\]

справедливе співвідношення

\[
H\left[y''\left(r\right)+\frac{1}{r}\ y'\left(r\right)\right]=-\xi^{2}H\left[y\right].
\]

Таким чином можна розв'язати спеціальне диференціальне рівняння з
змінними коефіцієнтами (рівняння Бесселя). 3. Перетворення Лапласа на
змінній t можна інтерпретувати як проекцію площини xt на осі x, що
призводить до У ц ь о м у

\pandocbounded{\includegraphics[keepaspectratio]{images/_page_104_Figure_9.jpeg}}

випадку початкові рівняння, граничні та початкові умови перетворюються
на нове диференціальне рівняння та нові граничні умови (див. діаграму
вище).\# ЗАВДАННЯ1. Переконайтеся, що наступна формула для перетворення
часткової похідної \(u_t\) є істинною:

\[
\mathscr{L}\left[u_{t}(x, t)\right] = sU(x, s) - u(x, 0).
\]

\begin{enumerate}
\def\labelenumi{\arabic{enumi}.}
\setcounter{enumi}{1}
\tightlist
\item
  Використати перетворення Лапласа для розв'язання задачі Коші
\end{enumerate}

\begin{verbatim}
\begin{aligned}
(\text{УЧП})\quad & u_t = \alpha^2 u_{xx}, && -\infty < x < \infty,\ 0 < t < \infty, \\
(\text{НУ})\quad & u(x,0) = \sin x, && -\infty < x < \infty.
\end{aligned}
\end{verbatim}

\begin{enumerate}
\def\labelenumi{\arabic{enumi}.}
\setcounter{enumi}{2}
\tightlist
\item
  За допомогою перетворення Лапласа за змінною t розв'язати задачу
\end{enumerate}

\begin{verbatim}
\begin{aligned}
(\text{УЧП})\quad & u_t = u_{xx}, && 0 < x < \infty,\ 0 < t < \infty, \\
(\text{ГУ})\quad & u(0,t) = \sin t, && 0 < t < \infty, \\
(\text{НУ})\quad & u(x,0) = 0, && 0 \leqslant x < \infty.
\end{aligned}
\end{verbatim}

Дайте фізичну інтерпретацію цієї задачі.4. Розв'язати задачу з граничним
значенням(ОДА)

\[
\frac{d^2U}{dx^2} - sU = A, \quad 0 < x < 1,
(\Gamma Y) \quad
\begin{cases} \frac{dU}{dx} (0) = 0 \\ U(1) = 0. \end{cases}
\]

Лекція 14\# ПРИНЦИП ДЮАМЕЛЯМЕТА ЛЕКЦІЇ: Показати, як переклад Лапласа
можна використовувати для ідентифікації характеристик диференціальних
рівнів. Зокрема, за допомогою алгебранічного перетворення образів
Лапласа, розвиток рівня часткового оберненого можна відкинути принципом
погляду Дюамеля. Інтерпретація цього принципу в теорії найвищих
диференціальних рівнів світу використовується для розробки певних
завдань УПП.технології.Давайте подивимось, як трансформація Лапласа
здійснюється спробою розвитку часткових диференціальних ліній, що
дозволяє нам розкрити фізичну сутність різних розвитків. За допомогою
трансформації Лапласа ми вже розробили принцип Дюамеля. Перше, що
потрібно зробити --- звернутися до принципу Дюамеля, ми бачимо проблему,
яка часто є проблемою\# Теплопровідність у стрижні, кінці якого
підтримуються при заданій температуріДуже часто необхідно знайти
температуру всередині області, якщо відомо, що на межі області вона
змінюється з часом відповідно до певного закону (граничні умови, залежні
від часу).На лівому кінці підтримується умова \(u(0,t)\)=0, а на правому
кінці температура задається функцією \(f(t)\).Рис. 14.1. Граничні умови,
що залежать від часу.Наприклад, розглянемо термоізольований стержень,
правий кінець якого підтримується у \(f(t)\) (див. рисунок 14.1).
Потрібно вирішити цю проблему(14.1)

\[
u_{t} = u_{xx}, \quad 0 < x < 1, \quad 0 < t < \infty,
(\Gamma Y) \quad \begin{cases} u(0, t) = 0, \\ u(1, t) = f(t), \end{cases} \quad 0 < t < \infty,
(HY) \quad u(x, 0) = 0, \quad 0 \le x \le 1.
\]

Сподіваємося, що розв'язання цієї задачі (14.1) буде легким, оскільки ми
знаємо розв'язок простішої задачі, де температура в кінці є
сталою:Насправді, якщо обидві задачі ((14.1) і (14.2)) розв'язуються
одночасно за допомогою перетворення Лапласа, то отримаємо визначний
результат: розв'язок задачі (14.1) виражається через розв'язання задачі
(14.2).Отже, розв'язуючи обидві задачі одночасно, ми отримуємо\# Легке
завдання (14.2) (Постійні солдати)Застосування перетворення -Лапласа до
задачі (14.2)

\[
\frac{d^2W}{dx^2} - sW(x) = 0,
W(0) = 0,
W(1) = 1/s
\]

.Пемеріль ODU

\[
W(x, s) = \frac{1}{s} \left[ \frac{\sinh(x \sqrt{s})}{\sinh(\sqrt{s})} \right].
\]

Пошук оберненого перетворення

\[
w(x, t) =
= x + \frac{2}{\pi} \sum_{n=1}^{\infty} \frac{(-1)^n}{n} e^{-(n\pi)^2 t} \times
\]

\(\times \sin(n\pi x)\) (розв'язання задачі з постійним ГІ).\# Важке
завдання (14.1) (Солдати, що залежать від часу)Застосування перетворення
Лапласа до задачі (14.1)

\[
\frac{d^2U}{dx^2} - sU(x) = 0,
U\left( \mathbf{0}\right) =0,
U(1) = F(s)
\]

.

\[
U(x, s) = F(s) \left[ \frac{\sinh(x \sqrt{s})}{\sinh(\sqrt{s})} \right].
\]

Розділіть і помножте на z

\[
U(x, s) = F(s) \left\{ s \left[ \frac{\sinh(x \sqrt{s})}{s \sinh(s)} \right] \right\}.
\]

Використовуючи співвідношення \(\mathscr{L}[w_t] = sW - w(x, 0),\)
отримуємо

\[
U(x, s) = F(s) \mathcal{L}[w_t].
\]

Відповідно,

\[
u(x, t) = \mathcal{L}^{-1} \{ F(s) \mathcal{L}[w_t] \} =
= \mathcal{L}^{-1} [ F(s) ] * \mathcal{L}^{-1}[w_t] =
= f(t) * w_t(t) =
= \int_0^t w_t(x, t-\tau) f(\tau) d\tau =
\]

(після інтегрування частинами)

\[
= \int_{0}^{t} w(x, t-\tau) f'(\tau) d\tau + f(0) w(x, t).
\]

(Розв'язок задачі з часово-залежними ГІ виражається через розв'язок
задачі з постійними ГІрозв'язано у задачах з константоюОтже, ми
висловили розв'язок задач із граничними умовами, що змінюються в часі,
через розв'язання задачі з постійними граничними умовами. Відповідні
формули такі:(14.3)плутанина у наступних формулах. Записати
диференціальне рівняння в похідних у нових координатах (~, т), мы
воспользуемся следующими формулами перехода;

Наприклад, якщо dy/dx = B/2A = 3, то y - 3x = c, а отже
\(\xi(x, y) = y - 3x\) . Ця функція задовольняє умову
\([\xi_x/\xi_y] = -3\) і, відповідно, повертає коефіцієнт
\(\overline{A}\) до нуля.Отже, половина роботи виконана: ми знайшли нову
координатну \(\xi = \xi(x, y)\) , яка забезпечила рівність
\(\overline{A} = 0\) . Залишається знайти \(\eta(x, y)\) таким чином,
щоб перетворити коефіцієнт \(\overline{B}\) на нуль.Час зупинитися тут.
Виявляється, що якщо вибрати \(B^2-4AC=0\) і \(\xi\) так, щоб
\(\overline{A}=0\) , то коефіцієнт \(\overline{B}\) автоматично стане
нулем 1). Тепер ми побачимо це.Коефіцієнт \(\overline{B}\) задається
виразом

\[
\overline{B}^2 - 4\overline{A}\overline{C} = (B^2 - 4AC) (\xi_x \eta_y - \xi_y \eta_x)^2 - \Pi \rho u u. \ \rho e \partial.
\overline{B} = 2 \left( V \overline{A} \xi_x + V \overline{C} \xi_y \right) \left( V \overline{A} \eta_x + V \overline{C} \eta_y \right).
\]

Функциональная диаграмма на рис. 15.5 делает эти формулы совсем
очевидными. Эта диаграмма особенно полезна при вычислении частных
производных функций и и и по переменным Р И С

\pandocbounded{\includegraphics[keepaspectratio]{images/_page_115_Picture_14.jpeg}}

\begin{verbatim}
15.5 Диаграмма, иллюстрирующая функциональную зависимость переменных  \xi   и   \tau  , поскольку   \xi   зависит и от x, и от t. Переменная   \tau   зависит только от t. Это все, что касается преобразований. Теперь подставим найденные выражения   u_t  ,   u_x   и   u_{vx}   в уравнение (15 2) и получим
\end{verbatim}

\[
\overline{B} = 2 \sqrt{\overline{A}} \left[ \sqrt{\overline{A}} \, \eta_x + \sqrt{\overline{C}} \, \eta_y \right].
\]

откуда

\[
u_{xy} + 2u_{xy} + u_{yy} = 0.
\]

\begin{verbatim}
.Следовательно, новая задача с начальным условием в персменных   \xi   и   \tau   имеет вид(УЧП)
\end{verbatim}

\[
\frac{dy}{dx} = -[\xi_x/\xi_y] = B/2A = 1.
\]

\begin{verbatim}
(НУ)   u(\xi, 0) = 1 - H(\xi), \quad -\infty < \xi < \infty.  Так как   \xi = x   при t = 0, то новые граничные условия совпадают со старыми. Эта задача была решена в лекции 12 с помощью интегрального преобразования Фурье, и ее решение записывается в виде
\end{verbatim}

\[
y = x + c
\]

\begin{verbatim}
где   \phi(\beta)   — начальное условие. Следовательно, в нашем случае
\end{verbatim}

\[
\eta = y
\]

После подстановки

\begin{verbatim}
\begin{cases}
 \xi = y - x, \\
 \eta = y
 \end{cases}
\end{verbatim}

получаем окончательный результат(15.3)

\[
\overline{C} = A \eta_x^2 + B \eta_x \eta_y + C \eta_y^2 = 1,
\overline{E} = A \eta_{xx} + B \eta_{xy} + C \eta_{yy} + D \eta_x + E \eta_y = 0,
\begin{cases} \frac{1}{2} \left[ 1 + \operatorname{erf} \left( \frac{-\xi}{2\sqrt{D\tau}} \right) \right], & \xi < 0, \\ \frac{1}{2} \operatorname{erfc} \left( \frac{\xi}{2\sqrt{D\tau}} \right), & 0 \leq \xi. \end{cases}
\overline{G} = G = 0.
\overline{A}u_{\xi\xi} + \overline{B}u_{\xi\eta} + \overline{C}u_{\eta\eta} + \overline{D}u_{\xi} + \overline{E}u_{\eta} + \overline{F}u = \overline{G}
\]

Графики этого решения для различных моментов времени изображены на рис.
15.6. Осталось сделать последний шаг --- записать решение нашей задачи в
координатах x и t. Подставляя

\[
u_{\eta} = f(\xi),
\]

\begin{verbatim}
,   \tau = t  в формулы (15.3), получаем(15.4)

$$
u(x, t) = \begin{cases} \frac{1}{2} \left[ 1 + \operatorname{erf}\left(\frac{Vt - x}{2\sqrt{Dt}}\right) \right], & Vt > x, \\ \frac{1}{2} \operatorname{erfc}\left(\frac{x - Vt}{2\sqrt{Dt}}\right), & Vt \leq x. \end{cases}
$$

Перед нами розв'язок задачі конвективної дифузії (15.2). Профіль концентрації рухається праворуч зі швидкістю $V$ і водночас розмивається дифузією з коефіцієнтом $D$.

#### ЗАВДАННЯ

1. Розв'яжіть задачу Коші

$$
\begin{array}{ll} u_t = u_{xx} - 2u_x, & -\infty < x < \infty, \ 0 < t < \infty, \\ u(x, 0) = \sin x, & -\infty < x < \infty. \end{array}
$$

2. Знайдіть розв'язок задачі Коші для рівняння конвективної дифузії

$$
\begin{array}{ll} u_t = u_{xx} - 2u_x, & -\infty < x < \infty, \ 0 < t < \infty, \\ u(x, 0) = e^x \sin x, & -\infty < x < \infty. \end{array}
$$

3. Знайдіть розв'язок задачі переносу, використовуючи перетворення координат з лекції 8.

## Приведення еліптичних рівнянь до канонічної форми

Розглянемо загальне рівняння

$$
Au_{xx} + Bu_{xy} + Cu_{yy} + Du_x + Eu_y + Fu = G
$$

але тепер у випадку $B^2 - 4AC < 0$. Переходом до нових незалежних змінних ми хочемо звести його до форми

$$
u_{\xi\xi} + u_{\eta\eta} = \varphi(\xi, \eta, u, u_{\xi}, u_{\eta}).
$$

Щоб знайти ці нові змінні, виконуємо ті самі обчислення, що й у двох попередніх випадках, отримуємо перетворене рівняння

$$
\overline{A}u_{\xi\xi} + \overline{B}u_{\xi\eta} + \overline{C}u_{\eta\eta} + \overline{D}u_{\xi} + \overline{E}u_{\eta} + \overline{F}u = \overline{G},
$$

і вимагаємо, щоб $\xi$ та $\eta$ задовольняли умовам $\overline{A} = \overline{C}$ і $\overline{B} = 0$. Безпосередньо знайти $\xi$ та $\eta$ тут складніше, тому шукаємо перетворення як композицію двох кроків.

# Перетворення 1

Спочатку введемо комплексні координати $\xi$ та $\eta$, щоб надати рівнянню вигляду

$$
u_{\xi\eta} = \psi(\xi, \eta, u, u_{\xi}, u_{\eta}).
$$

Для цього розв'язуємо характеристичні рівняння

$$
\frac{dy}{dx} = \frac{B - \sqrt{B^2 - 4AC}}{2A},
$$

де $B^2 - 4AC < 0$, і отримуємо

$$
\xi(x, y) = \text{const}, \qquad \eta(x, y) = \text{const}.
$$

# Перетворення 2

Далі переходимо від $(\xi,\eta)$ до $(\alpha,\beta)$ за формулами

$$
\alpha = \frac{\xi + \eta}{2}, \qquad \beta = \frac{\xi - \eta}{2i}.
$$

У результаті отримуємо канонічну форму еліптичного рівняння.

# Зведення рівняння y^2u_{xx} + x^2u_{yy} = 0 до його канонічної формиРозглянемо рівняння
\end{verbatim}

у яких \(A=y^2\) , B=0, \(C=x^2\) , D=E=F=G=0. Дискримінантний
\(B^2-4AC\) дорівнює \(-4x^2y^2\) , і ми зведемо це рівняння до
канонічної форми в першому квадранті x\textgreater0, y\textgreater0.КРОК
1 (Перше перетворення)Запишемо характеристичні рівняння

\[
\frac{dy}{dx} = \frac{B + \sqrt{B^2 - 4AC}}{2A} = \frac{\sqrt{-4x^2y^2}}{2y^2} = i\frac{x}{y}
\]

Замініть змінні \(\xi = x - Vt\) , \(\tau = t\) . У нашому випадку
інтеграл можна обчислити явно

\[
y^2 + ix^2 = \text{const},
\]

Такий інтеграл уже зустрічався у лекції 12. Можливо, читачу буде
зручніше звернутися до таблиць перетворення Фур'є.\# ОДНОВИМІРНЕ
ХВИЛЬОВЕ РІВНЯННЯ (ГІПЕРБОЛІЧНІ РІВНЯННЯ) МЕТАЛЕКЦІЇ: Показати, що
коливання навантаження описуються рівнянням

Від

і що це рівняння є наслідком законів Ньютона. Також розглядаються інші
типи хвильового рівняння:

\begin{verbatim}
\eta(x, y) = y^2 - ix^2. (Неважливо, яка з двох функцій називається \xi , а яка — \eta ; ви можете поміняти їх місцями, якщо хочете.) Після цього перетворення початкове рівняння набуває вигляду
\end{verbatim}

Ми не звертатимемо уваги на це рівняння (яке є комплексним гіперболічним
рівнянням), а одразу перейдемо до другого перетворення.КРОК 2. (Друга
трансформація)При здійсненні другої трансформації ми маємо

\[
\begin{split} u_{tt} &= \alpha^2 u_{xx} + F\left(x, \ t\right), \\ u_{tt} &= \alpha^2 u_{xx} - \beta u_t - \lambda u + F\left(x, \ t\right). \end{split}
\]

(реальна частина \(\xi\) і \(\eta\) ),
\(\beta = \frac{\xi - \eta}{2\iota} = x^2\) (уявна частина \(\xi\) і
\(\eta\) ).Щоб спростити нотацію, найкраще перепозначити
\((\alpha,\beta)\) \((\xi,\eta)\) окремо і записати отримане
\(\epsilon\) перетворення як

\begin{verbatim}
КРОК 3 (Отримання нового рівняння)Нову канонічну форму можна отримати, обчисливши всі коефіцієнти \overline{A}, \overline{B}, \overline{C}, \overline{D}, \overline{E}, \overline{F} і \overline{G} у рівнянні
\end{verbatim}

До цього часу ми розглядали фізичні процеси, які описуються одномірними
параболічними рівняннями (задачі дифузії). Тепер ми почнемо вивчати
інший клас диференціальних рівнянь у похідних --- гіперболічні рівняння.
Почнемо з одновимірного хвильового рівняння, яке описує, зокрема,
поперечні коливання струни.\# Вібрації струнРозглянемо невеликі
коливання струни з нерухомими кінцями. Припустимо, що струна щільно
натягнута, зроблена з однорідного матеріалу і коливається в одній
площині (рис. 16.1).

\pandocbounded{\includegraphics[keepaspectratio]{images/_page_119_Figure_10.jpeg}}

Фіг. 16.1. Поперечні коливання струни: a --- зміщення струни з положення
рівноваги.Ігноруємо вплив гравітації на струну. Щоб побудувати
математичну модель коливань струн, розглянемо всі сили, що діють на
невелику частину струни (рис. 16.2). Виявляється, що хвильове рівняння
--- це не більше ніж закон руху Ньютона (зміна імпульсу дорівнює сумі
діючих сил), застосований до струни. Дивлюся на рис. На рисунку 16-1 ми
можемо представити сили, що діють на струну, у напрямках,
перпендикулярних осі x.

\pandocbounded{\includegraphics[keepaspectratio]{images/_page_120_Figure_2.jpeg}}

Фіг. 16.2. Невелика частина \((x, x + \Delta x)\) осцилюючою струною.1.
Сила, що виникає через натяг струни ( \(\alpha^2 u_{xx}\) ). Значення
поперечної компоненти сили натягу визначається відповідно до формулиСила
=

\begin{verbatim}
за формулами (41.5) з \xi = y^2 та \eta = x^2 . Зробивши це, отримуємо канонічну форму еліптичного рівняння
\end{verbatim}

\(\simeq T [u_x(x + \Delta x, t) - u_x(x, t)].\) 2. Внешняя сила
\(F (x, t)\).Зовнішня сила \(F(x, t)\), прикладена до струни, може
довольно залежати від x і t. Нижче наведено деякі приклади зовнішніх
сил:а) гравитационная сила \(F(x, t)\) = -mg;b) імпульси, розподілені
вздовж струни та діючи на неї в різний час:c) сили, що виникають
внаслідок дії звукової хвилі на шкіру барабана (пізніше ми детально
вивчаємо коливання мембрани, які описуються хвильовим рівнянням).3. Сили
тертя, що діють на струну (--- \(\beta u_t\) ). Якщо струна вібрує у
середовища, виникає сила опору, яка пропорційна швидкості \(u_t\) .4.
Возвращающая сила (- үи).Ця сила спрямована протилежно зміщенню струни.
Якщо зміщення u додатне, то сила від'ємна.Якщо тепер застосувати
рівняння руху Ньютона до невеликої ділянки струни, отримаємо

\section{\texorpdfstring{ПРИМІТКИ1. Аналогічно можна класифікувати
лінійні рівняння другого порядку з \emph{трьох або більше змінних}, але
для цього потрібен матричний аналіз. Наприклад, рівняння другого порядку
з трьома
змінними}{ПРИМІТКИ1. Аналогічно можна класифікувати лінійні рівняння другого порядку з трьох або більше змінних, але для цього потрібен матричний аналіз. Наприклад, рівняння другого порядку з трьома змінними}}\label{ux43fux440ux438ux43cux456ux442ux43aux4381.-ux430ux43dux430ux43bux43eux433ux456ux447ux43dux43e-ux43cux43eux436ux43dux430-ux43aux43bux430ux441ux438ux444ux456ux43aux443ux432ux430ux442ux438-ux43bux456ux43dux456ux439ux43dux456-ux440ux456ux432ux43dux44fux43dux43dux44f-ux434ux440ux443ux433ux43eux433ux43e-ux43fux43eux440ux44fux434ux43aux443-ux437-ux442ux440ux44cux43eux445-ux430ux431ux43e-ux431ux456ux43bux44cux448ux435-ux437ux43cux456ux43dux43dux438ux445-ux430ux43bux435-ux434ux43bux44f-ux446ux44cux43eux433ux43e-ux43fux43eux442ux440ux456ux431ux435ux43d-ux43cux430ux442ux440ux438ux447ux43dux438ux439-ux430ux43dux430ux43bux456ux437.-ux43dux430ux43fux440ux438ux43aux43bux430ux434-ux440ux456ux432ux43dux44fux43dux43dux44f-ux434ux440ux443ux433ux43eux433ux43e-ux43fux43eux440ux44fux434ux43aux443-ux437-ux442ux440ux44cux43eux43cux430-ux437ux43cux456ux43dux43dux438ux43cux438}

де \(\rho\) --- лінійна густина струни. Поділимо обидві частини цього
співвідношення на \(\Delta x\) і поставимо \(\Delta x\) --- 0.
Результатом є добре відоме телеграфне рівняння(16.1)

згідно з цією класифікацією, буде параболічним, а рівняння

У цьому рівнянні слід поділити \(\alpha\) , \(\beta\) , \(\gamma\) і
\(F(x,t)\) на \(\rho\) , але для простоти залишимо старі позначки для
цих величин. Отримано бажане рівняння. Ось інтерпретація найпростішого
хвильового рівняння

\begin{verbatim}
- Гіперболічно.2. Наш інтерес до класифікації диференціальних рівнянь у похідних пов'язаний із тим, що три основні типи рівнянь описують три різні типи фізичних явищ, і ми хотіли б отримати математичну класифікацію цих трьох класів фізичних явищ.# ЗАВДАННЯ1. Які з цих параболічних і еліптичних рівнянь записані у канонічній формі?a) u_t = u_{xx} + hu,б) u_{xy} + u_{xx} + 3u = \sin x,в) u_{xx} + 2u_{yy} = 0,г) u_{xx} = \sin^2 x?2. Перетворіть параболічне рівняння u_{xx} + 2u_{xy} + u_{yy} + u = 2 до канонічної форми.3. Перетворіть еліптичне рівняння u_{xx} + 2u_{yy} + x^2 u_x = e^{-x^2/2} до канонічної форми.# МЕТОД МОНТЕ-КАРЛО (ВСТУП)МЕТА ЛЕКЦІЇ: Пояснити основні ідени, що лежать в основі методу Монте-Карло, і показати, як цей метод можна використовувати для розв'язання різних задач. Виявляється, що за допомогою азартних ігор можна змоделювати (загалом, на комп'ютері) наближені розв'язки реальних задач. Як простий приклад можна вказати обчислення інтеграла
\end{verbatim}

основанную на интуитивных соображениях.\# Інтуїтивна інтерпретація
хвильового рівнянняЧитач може запитати: чому рівняння типу (16.2) має
описувати щось подібне до вібрацій скрипкової струни? Щоб зрозуміти це,
необхідно зрозуміти Ф і

\pandocbounded{\includegraphics[keepaspectratio]{images/_page_121_Picture_9.jpeg}}

г. 16 3. Інтерпретація хвильового рівняння \(u_{tt} = \alpha^2 u_{xx}\)
\(u_{tt}\) Equation (16.2) можна інтерпретувати так: прискорення струни,
викликане натягом, у кожній точці більше, чим більша увігнутість струни
\(u_{xx}\) у заданій точці (коефіцієнт пропорційності сталий
\(\alpha^2 = T/\rho\) ) (див. рис. 133). \# ПОВАГА 1. Хвильове рівняння
\(u_{tt} = \alpha^2 u_{xx}\) також описує поздовжні та крутильні
коливання стрижня. У поздовжніх осцилях зміщення паралельні елементу, і
значення поздовжнього зміщення відносно положення рівноваги позначається
через \(u\left(x,\,t\right)\) . Такі коливання виникають, наприклад,
коли молоток б'ють по кінці стрижня. Торсіональнийные колебания

\begin{verbatim}
кидаючи дротик у одиничний квадрат \{(x, y): 0 < x < 1, 0 < y < 1\} Після приблизно 100 кидків відносна частина дротиків, що опускається нижче кривої y = x^2 , вважається приблизним значенням інтеграла. Узагальнення такої гри (з випадковим киданням дрогіка) завжди містить процедуру генерації послідовності випадкових чисел.У цій лекції ми коротко опишемо метод Монте-Карло, а в наступній лекції покажемо, як його використовувати для розв'язання диференціальних рівнянь у похіднихМетоди Монте-Карло зазвичай називають групою методів розв'язання детермінованих задач (тобто задач без випадковості), у яких суттєво використовуються елементи випадковості. Загальна філософія етнічних методів показана на рис. 42 1.Ймовірнісна граДеторміноване завдання \rho — розв'язання задачі (наприклад, обчислення інтеграла, розв'язання UBP тощо).Результат = \hat{p} Наближення p Рішення = p Рис. 42.1. Загальна теорія методу Монте-Карло.# Обчислення інтегралаПрипустимо, ми хочемо обчислити інтеграл
\end{verbatim}

Тут \emph{k} --- це модуль Юнга, який характеризує еластичність
матеріалу. Матеріали з великим модулем Юнга коливаються на вищих
частотах. Звукові хвилі --- це поздовжні хвилі.2. Якщо лінійна густина
струни \(\rho(x)\) залежить від координати, то хвильова аугура
записується як

\begin{verbatim}
(детерміноване завдання). Щоб використати метод Монте-Карло, потрібно винайти гру на удачу, яка дає приблизне значення цього інтеграла.Звісно, можна придумати багато різних ігор. Остаточна версія залежить від точності розрахунків, простоти гри тощо. Очевидний варіант гри для обчислення інтегралу — кинути дротик у прямокутний R = \{(x,y): a \le x \le b, 0 \le y \le \max f(x)\} (див. рис. 42.2)
\end{verbatim}

\pandocbounded{\includegraphics[keepaspectratio]{images/_page_126_Figure_14.jpeg}}

. 17.1. Волна \(u(x, t)\) = \(e^{-(x-ct)^2}\) движется слева
направопроизвольной формы, движущихся в прогивоположных направлениях со
скоростью c.~Например, функции

\[
= \int_{0}^{d} \left[\frac{\partial F}{\partial \overline{y}} \eta\left(x\right) + \frac{\partial F}{\partial \overline{y'}} \eta'\left(x\right)\right] dx = 0.
\]

(волна, Лвижущаяся слева направо),

\[
\frac{d\varphi(e)}{de} \equiv \int_{a}^{b} \left\{ \frac{\partial F}{\partial \overline{y}} - \frac{\partial}{\partial x} \left[ \frac{\partial F}{\partial \overline{y'}} \right] \right\} \, \eta(x) \, dx = 0.
\]

(волна, движущаяся справа налево),

\[
\frac{\partial F}{\partial \overline{y}} - \frac{d}{dx} \left[ \frac{\partial F}{\partial \overline{y'}} \right] = 0
\]

(две волны, движущиеся в противоположных направлениях)являются типичными
решениями волиового уравнения. На рисунке 17.1 изображена простейшая
бегущая волна.ШАГ 4 (Учет начальных условий).Напомним читателю общий
метод решения задачи Коши в теории обыкновенных дифференциальных
уравнений. Сначала находят общее решение уравнения, а затем оно
подставляется в начальные условия для того, чтобы найти конкретные
значения произвольных постоянных. У нас ситуация аналогична. Для решения
задачи Коши для волнового уравнения мы подставим общее решение

\[
\frac{\partial F}{\partial y} - \frac{d}{dx} \left[ \frac{\partial F}{\partial y'} \right] = 0
\]

(содержащее две произвольные функции) в начальные условия

\[
J[y] = \int_{0}^{1} [y^{2} + y'^{2}] dx.
\]

\(u_1\)(x, 0) = \(g(x)\), чтобы найти конкретные выражения для
произвольных функций ф и ф. После подстановки получаем(17.5)

\[
F_y - \frac{d}{dx} F_{y'} = 0,
\]

\begin{verbatim}
Если теперь проинтегрировать второе уравнение (17.5), то получится алгебранческая система двух уравнений относительно неизвестных функций   \varphi(x)   и   \psi(x)  . В самом деле, после интегрирования в пределах от   x_0   до x получаем(17.6)
\end{verbatim}

\[
\begin{aligned} F_y &= 2y, \ F_{y'} &= 2y'. \end{aligned}
\]

\begin{verbatim}
Решаем теперь уравнение (17.6) совместно с первым уравнением (17.5) и получаем следующие выражения для функций   \varphi(x)   и   \psi(x)  :(17.7)
\end{verbatim}

\[
2y - \frac{d}{dx}(2y') = 0,
y'' - y = 0.
\]

Отже, загальне розв'язання задачі (17.1) задається формулою(17.8)

\[
\overline{y}(x) = 0.42e^x - 0.42e^{-x}
\]

Мета досягнута. Ми отримали загальне розв'язання Вадачі (17.1).
Розв'язок у вигляді (17.8) зазвичай називають формулою д'Аламбера.
Читача запрошують самостійно слідкувати за розташуванням меж
інтегрування, щоб переконатися, що вони змінюються з x-ct на x+ct.
Проблема повністю вирішена.Перед завершенням лекції ми наведемо кілька
прикладів, які показують, як застосовувати формулу д'Аламбера до
конкретних валахів.\# Приклади застосувань формули д'Аламбера\# 1. Рух
синузондалних хвильРассмотрим начальные условия вида

\[
J[y] = \int_a^b y \sqrt{1 + y'^2} \, dx,
\]

\(u_1(x, 0) = 0.\) По формуле Даламбера получаем решение

\[
\overline{y}(x) = \alpha \operatorname{ch}[(x-\beta)/\alpha],
\]

Її можна інтерпретувати так. Початкове переміщення струни
\(u(x, 0) = \sin x\) поділяється на дві однакові частини

\[
\overline{y}(x) = 0.42e^x - 0.42e^{-x}
\]

H \(\frac{\sin x}{2}\) .Кожна з частин поширюється зі швидкістю c у
вигляді хвиль, що рухаються. Одна з хвиль рухається зліва направо, а
інша --- у протилежному напрямку. Читача запрошують замислитися, як
виглядає отримана хвиля у цьому випадку.\# 2. Распространение
простейшего прямоугольного импульсаВ этом случае начальные условия
задаются следующим образом:

\[
J[y] = \int_{0}^{1} [y^{2} + y'^{2}] dx.
\]

\(u_t(x, 0) = 0\) (в остальных точках).Оскільки початкове зміщення
розділяється на дві півхвилі, що рухаються в протилежних напрямках,
отриманий рух буде таким, як показано на рис. 17.2.

\pandocbounded{\includegraphics[keepaspectratio]{images/_page_129_Figure_2.jpeg}}

Фіг. 17.2. Початкове збурення розділяється на дві хвилі, що рухаються.\#
3. Задана начальная скоростьТепер припустимо, що в початковий момент
струна перебуває в рівноваговій позиції. Надамо струні початкову
швидкість (як у фортепіано) форми \(\sin x\) , тобто

\[
\int\limits_{t_1}^{t_2}
J\left[u\right] = \iint\limits_{D} F\left(x,\ y,\ u,\ u_{x},\ u_{y}\right) \, dx \, dy.
\]

Решение находим по формуле Даламбера

\[
F_{u} - \frac{\partial}{\partial x} F_{ux} - \frac{\partial}{\partial y} F_{uy} = 0.
\]

Оно представляет сумму двух бегущих косинусоидальных воли. Читатель
должен спросить себя, кажется ли ему это решение разумным.На этом лекция
17 заканчивается. В следующей лекции мы покажем, как формула Даламбера
помогает дать интерпретацию в плоскости переменных x и t.

\section{ЗАУВАЖЕННЯ - 1. Общее решение уравнения с частными производными
второго порядка содержит две произвольные функции, а общее решение
обыкновенного дифференциального уравнения содержит две произвольные
постоянные. Значит, у уравнения с частными производными больше решений,
чем у обыкновенного дифференциального
уравнения.}\label{ux437ux430ux443ux432ux430ux436ux435ux43dux43dux44f---1.-ux43eux431ux449ux435ux435-ux440ux435ux448ux435ux43dux438ux435-ux443ux440ux430ux432ux43dux435ux43dux438ux44f-ux441-ux447ux430ux441ux442ux43dux44bux43cux438-ux43fux440ux43eux438ux437ux432ux43eux434ux43dux44bux43cux438-ux432ux442ux43eux440ux43eux433ux43e-ux43fux43eux440ux44fux434ux43aux430-ux441ux43eux434ux435ux440ux436ux438ux442-ux434ux432ux435-ux43fux440ux43eux438ux437ux432ux43eux43bux44cux43dux44bux435-ux444ux443ux43dux43aux446ux438ux438-ux430-ux43eux431ux449ux435ux435-ux440ux435ux448ux435ux43dux438ux435-ux43eux431ux44bux43aux43dux43eux432ux435ux43dux43dux43eux433ux43e-ux434ux438ux444ux444ux435ux440ux435ux43dux446ux438ux430ux43bux44cux43dux43eux433ux43e-ux443ux440ux430ux432ux43dux435ux43dux438ux44f-ux441ux43eux434ux435ux440ux436ux438ux442-ux434ux432ux435-ux43fux440ux43eux438ux437ux432ux43eux43bux44cux43dux44bux435-ux43fux43eux441ux442ux43eux44fux43dux43dux44bux435.-ux437ux43dux430ux447ux438ux442-ux443-ux443ux440ux430ux432ux43dux435ux43dux438ux44f-ux441-ux447ux430ux441ux442ux43dux44bux43cux438-ux43fux440ux43eux438ux437ux432ux43eux434ux43dux44bux43cux438-ux431ux43eux43bux44cux448ux435-ux440ux435ux448ux435ux43dux438ux439-ux447ux435ux43c-ux443-ux43eux431ux44bux43aux43dux43eux432ux435ux43dux43dux43eux433ux43e-ux434ux438ux444ux444ux435ux440ux435ux43dux446ux438ux430ux43bux44cux43dux43eux433ux43e-ux443ux440ux430ux432ux43dux435ux43dux438ux44f.}

\begin{itemize}
\tightlist
\item
  2 Замена переменных для приведения уравнения с частными производными к
  более простому виду является одним из общих методов теории. Новые
  координаты (ξ, η) принято называть каноническими координатами. Мы
  будем обращаться к ним и в дальнейшем, особенно при изучении
  гиперболических уравнений.
\end{itemize}

\begin{verbatim}
- 3 Метод нахождения общего решения уравнения с частными производными и подстановки этого решения в начальные условия не является общим методом решения уравнений с частными производными. Решение, которое рассматривается в данной лекции,— единственный пример использования этого метода. Обычно невозможно найти общее решение уравнения с частными производными, но даже если удается это сделать, оказывается очень сложным подставить это решение в краевые условия.# ЗАДАЧН1. Проверьте, что решение Даламбера (17-8) удовлетворяет начальным условиям задачи (17.1).2. Подставьте (17.7) в общее решение   u(x, t) = \psi(x - ct) + \psi(x + \epsilon t)   и получите формулу Даламбера.3. Найдите решение задачи Коши:
\end{verbatim}

\[
J[y] = \int_{0}^{1} \sqrt{1 + y'^{2}} dx.
\]

Как выглядит это решение в различные моменты времени? 4. Найдите решение
задачи Коши:(УЧП)

\[
\int_{t_{1}}^{t_{2}} [KE - PE] dt = \frac{1}{2} \int_{t_{1}}^{t_{2}} [m\dot{y}^{2} - ky^{2}] dt
\]

(НУ)

\[
J[y] = \int_{0}^{\pi/2} [y'^{2} - y^{2}] dx
\]

Постройте график решения в различные моменты времени. 5. Решая совместно
уравнения (17.6) и первое уравнение из (17.5),получите формулы (17.7).\#
ФОРМУЛА ДАЛАМБЕРА (ПРОДОЛЖЕНИЕ)ЦЕЛЬ ЛЕКЦИИ! Показать, как с помощью
формулы Даламбера можно найти решение волнового уравнения на
полуограниченной прямой(УЧП)

\[
J[u] = \iint_D F[x, y, u, u_x, u_y) \, dx \, dy,
\]

\begin{verbatim}
(ГУ)   u(0, t) = 0, \quad 0 < t < \infty,
\end{verbatim}

(НУ)

\[
F_{u} - \frac{\partial}{\partial x} F_{uv} - \frac{\partial}{\partial y} F_{uv} = 0.
\]

Дополнительно дается интерпретация формулы Даламбера в координатной
плоскости x, t.На прошлой лекции было показано, что если заданы
начальное смещение струны \(u(x, 0)\) = \(f(x)\) и начальная скорость
\(u_t\)(x, 0) = \(g(x)\) , то выражение(18.i)

\[
u_{xx} + u_{yy} = f
\]

описывает смещение струны в произвольный момент времени. В этой лекции
мы познакомим читателя с одной интересной интерпретацией этой формулы в
плоскости переменных x и t (пространственновременная или фазовая
плоскость) и покажем, как изменить формулу Даламбера, чтобы получить
решение задачи о колебаниях полубесконечной струны.Начнем с
интерпретации формулы (18.1) в плоскости переменных x и t.\#
Пространственно-временная интерпретация формулы ДаламбераНа предыдущей
лекции мы показали, что решение задачи Коши(18.2)

\[
J[u] = \iint_{D} [u_{x}^{2} + u_{y}^{2} + 2uf] dx dy.
J[u] = \int_{0}^{1} \int_{0}^{1} \left[ u_{1}^{2} + u_{y}^{2} \right] dx dy
u(x, t) = \frac{1}{2} [f(x-ct) + f(x+ct)] + \frac{1}{2c} \int_{x-ct}^{x+ct} g(\xi) d\xi.
u_{xx} + u_{yy} = 0
u(x, 0) = f(x),
u_{xx} + u_{yy} = 0
u(x, t) = \frac{1}{2} [f(x-ct) + f(x+ct)].
J[u] = \int_{0}^{1} \int_{0}^{1} \left[ u_{x}^{2} + u_{y}^{2} \right] dx dy
x-ct=x_0-ct_0
\]

(PPP)

\[
(Y \text{ ЧП}) \qquad u_{tt} = c^2 u_{xx}, \quad -\infty < x < \infty, \quad 0 < t < \infty,
J[u] = \iint_{D} [u] + u_{\eta} + 2uf ] dx dy.
(HY) \qquad
u_{xx} + u_{yy} = f
J[u] = \int_{0}^{1} \int_{0}^{1} [u' + u' + 2uf] dx dy,
J[u] = \int_{0}^{1} \int_{0}^{1} \left[ u_{x}^{2} + u_{y}^{2} + 2uf \right] dx dy,
\]

Фіг. 18.2 Розв'язок зображено на площині змінних x, t, а не u, x, як у
попередній лекції.

\pandocbounded{\includegraphics[keepaspectratio]{images/_page_133_Figure_5.jpeg}}

Фіг. 18.2. Розв'язок задачі Коші (18.2) у xt-площині a --- це передній
край хвилі, що рухається вліво; b --- це задній край хвилі; b --- це
передній край хвилі, що рухається вправоТепер інтерпретуємо формулу
д'Аламбера, де початкове зміщення дорівнює нулю, а початкова швидкість
довільна.СЛУЧАЙ 2 (Начальное смещение равно нулю, начальная скорость
произвольна).Рассмотрим начальные условия вида

\[
u_n(x, y) = a_1 \varphi_1(x, y) + a_2 \varphi_2(x, y) + \ldots + a_n \varphi_n(x, y),
\]

\(u_t(x, 0) = g(x).\) Решение в этом случае имеет вид

\[
\varphi_1(x, y) = xy(1-x)(1-y)
\]

Значення значення значення U в точці
\((x_{\mathrm{ }e},\ t_{\mathrm{ }e})\) можна інтерпретувати як інтеграл
початкової швидкості від \(x_{\mathrm{ }e}-ct_{\mathrm{ }e}\) до P і C

\pandocbounded{\includegraphics[keepaspectratio]{images/_page_134_Figure_2.jpeg}}

. 18.3. Интерпретация задачи с начальной скоростью в x1-плоскости.
\(x_{\bullet}+ct_{\bullet}\) (рис. 18.3). Наприклад, рис. 18.4 надає
графічне тлумачення розв'язку задачі(13.4)

\[
u_{1}(x, y) = a_{1}xy(1-x)(1-y),
u_{3}(x, y) = xy(1-x)(1-y)[a_{1} + a_{2}x + a_{3}y],
\vdots
u(x, t) = \frac{1}{2c} \int_{x-ct}^{x+ct} g(\xi) d\xi =
J[u_n] = \int_0^1 \int_0^1 \left\{ \left[ \sum_{j=1}^n a_j \frac{\partial \varphi_j}{\partial x} \right]^2 + \left[ \sum_{j=1}^n a_j \frac{\partial \varphi_j}{\partial y} \right]^2 + 2f \sum_{j=1}^n a_j \varphi_j \right\} dx dy
= \frac{1}{2c} \int_{x-ct}^{x+ct} 0 d\xi, \quad (x, t) \in \text{ области 1,}
\frac{\partial J\left[u_{n}\right]}{\partial a_{1}} = 2 \int_{0}^{1} \int_{0}^{1} \left\{ \sum_{j=1}^{n} \left[ \frac{\partial \varphi_{j}}{\partial x} \frac{\partial \varphi_{j}}{\partial x} + \frac{\partial \varphi_{j}}{\partial y} \frac{\partial \varphi_{1}}{\partial y} \frac{\partial \}{\partial y} \right] a_{j} + f \varphi_{1} \right\} dx dy = 0,
= \frac{1}{2c} \int_{x-ct}^{x+ct} d\xi, \quad (x, t) \in \text{ области 2,}
\]

Все это представляется изрядно сложным, но если переписать все эти
уравнения в матричной форме, то получится система линейных уравнений

\[
= \frac{1}{2c} \int_{x-ct}^{-1} d\xi, \quad (x, t) \in \text{ области 3,}
\]

\begin{verbatim}
где  A = (A_{IJ})  матрица размером  n \times n , элементы которой вычисляются по формулам(45.3)
\end{verbatim}

\[
= \frac{1}{2c} \int_{x-ct}^{1} d\xi, \quad (x, t) \in \text{ области 4,}
\]

b = (\(b_t\)) --- вектор с компонентами

\[
= \frac{1}{2c} \int_{x-ct}^{x+ct} 0 d\xi, \quad (x, t) \in \text{ области 5,}
\]

\begin{verbatim}
а  a=(a_i) — неизвестный вектор, компоненты которого являются коэффициентами в приближенном решении  a_n(x,y)=a_1\varphi_1(x,y)+a_2\varphi_2(x,y)+\ldots+a_n\varphi_n(x,y) .ШАГ 3. Решаем линейную систему Aa = b относительно коэффициентов  a_1, a_2, \ldots, a_n  и получаем приближенную минимизирующую функцию
\end{verbatim}

\[
= \frac{1}{2c} \int_{x-ct}^{x+ct} d\xi, \quad (x, t) \in \text{ области 6.}
\]

а эначит, и приближенное решение задачи Дирихле.

\section{ЗАУВАЖЕННЯ}\label{ux437ux430ux443ux432ux430ux436ux435ux43dux43dux44f}

\begin{enumerate}
\def\labelenumi{\arabic{enumi}.}
\tightlist
\item
  В этой лекции методом Ритца минимизировался функционал, вависящий от
  функции двух переменных. Точно так же можно минимизировать функционалы
  вида
\end{enumerate}

\[
y(0) = 0, \quad y(1) = 1.
\begin{array}{llll} (\text{Y}4\Pi) & u_{tt} = c^2 u_{xx}, & 0 < x < \infty, & 0 < t < \infty, \\ (\Gamma \text{Y}) & u(0, t) = 0, & 0 < t < \infty, \\ (\text{H}\text{Y}) & \begin{cases} u(x, 0) = f(x), \\ u_t(x, 0) = g(x). \end{cases} & 0 < x < \infty, \\ \end{array}
\varphi_i(0) = 0,
J[u] = \int_{0}^{1} \int_{0}^{1} [u^{2} + u^{2}_{y}] dx dy,
\]

\section{Розв'язання задачі для напівнескінченного рядка за допомогою
формули д'АламбераНаше завдання --- знайти хвильові рухи
напівнескінченної струни, лівий кінець якої фіксований, за заданих
початкових умов. Для розв'язання задачі (18.6) ми використаємо той самий
метод, що й у пошуку загального розв'язку хвильового рівняння на всій
лінії. Заміна загального
розв'язку}\label{ux440ux43eux437ux432ux44fux437ux430ux43dux43dux44f-ux437ux430ux434ux430ux447ux456-ux434ux43bux44f-ux43dux430ux43fux456ux432ux43dux435ux441ux43aux456ux43dux447ux435ux43dux43dux43eux433ux43e-ux440ux44fux434ux43aux430-ux437ux430-ux434ux43eux43fux43eux43cux43eux433ux43eux44e-ux444ux43eux440ux43cux443ux43bux438-ux434ux430ux43bux430ux43cux431ux435ux440ux430ux43dux430ux448ux435-ux437ux430ux432ux434ux430ux43dux43dux44f-ux437ux43dux430ux439ux442ux438-ux445ux432ux438ux43bux44cux43eux432ux456-ux440ux443ux445ux438-ux43dux430ux43fux456ux432ux43dux435ux441ux43aux456ux43dux447ux435ux43dux43dux43eux457-ux441ux442ux440ux443ux43dux438-ux43bux456ux432ux438ux439-ux43aux456ux43dux435ux446ux44c-ux44fux43aux43eux457-ux444ux456ux43aux441ux43eux432ux430ux43dux438ux439-ux437ux430-ux437ux430ux434ux430ux43dux438ux445-ux43fux43eux447ux430ux442ux43aux43eux432ux438ux445-ux443ux43cux43eux432.-ux434ux43bux44f-ux440ux43eux437ux432ux44fux437ux430ux43dux43dux44f-ux437ux430ux434ux430ux447ux456-18.6-ux43cux438-ux432ux438ux43aux43eux440ux438ux441ux442ux430ux454ux43cux43e-ux442ux43eux439-ux441ux430ux43cux438ux439-ux43cux435ux442ux43eux434-ux449ux43e-ux439-ux443-ux43fux43eux448ux443ux43aux443-ux437ux430ux433ux430ux43bux44cux43dux43eux433ux43e-ux440ux43eux437ux432ux44fux437ux43aux443-ux445ux432ux438ux43bux44cux43eux432ux43eux433ux43e-ux440ux456ux432ux43dux44fux43dux43dux44f-ux43dux430-ux432ux441ux456ux439-ux43bux456ux43dux456ux457.-ux437ux430ux43cux456ux43dux430-ux437ux430ux433ux430ux43bux44cux43dux43eux433ux43e-ux440ux43eux437ux432ux44fux437ux43aux443}

\[
u_{xx} + u_{yy} = 1
\]

в начальные условия, находим (как в предыд, щей лекции)(18.7)

\[
J[y] = \int_{0}^{1} [y^{2} + {y'}^{2}] dx, \ y(0) = 0, \ y(1) = 1,
z(x) = (1 - x) y(x),
\]

\begin{verbatim}
"Теперь перед нами проблема, которой не было в задаче о бесконечной струне. Решение   u\left(x,\t\right)   должно быть определено всюду только внутри первого квадранта   (x>0,\,t>0)   плоскости переменных x и t. Значит, мы должны уметь вычислить функции
\end{verbatim}

\[
u_{xx} + u_{yy} = 0
\]

\begin{verbatim}
при всех   -\infty < x-ct < \infty  ,   \psi(x+ct)   при всех   0 < x+ct < \infty  .К сожалению, первая из формул (18.7) позволяет вычислять   \varphi(x-ct)   только при   x-ct\geqslant 0  , поскольку в начальных данных функции f(x) и g(x) определены только для положительных значений аргумента.Если   x-ct \geqslant 0  , то после подстановки (18.7) в общее решение
\end{verbatim}

\[
J[u] = \int_{0}^{1} \int_{0}^{1} [u_{x}^{2} + u_{y}^{2}] dx dx.
\]

получаем

\[
u_{xx} + u_{yy} = \sin(\pi x)
u(x, y) = \left[Ae^{\pi y} + Be^{-\pi y} - \frac{1}{n^2}\right] \sin(\pi x),
\begin{split} u\left(x,\,t\right) &= \phi\left(x-ct\right) + \psi\left(x+ct\right) = \\ &= \frac{1}{2} \left[ f\left(x-ct\right) + f\left(x+ct\right) \right] + \frac{1}{2c} \int\limits_{x-ct}^{x+ct} g\left(\xi\right) \, d\xi. \end{split}
\]

\Дельта u + u\^{}2 = 0

\[
\Delta u = 0
\]

\begin{verbatim}
Что же делать, если x < ct? Пришла пора воспользоваться граничным условием u(0,t)=0. С помощью этого граничного условия мы доопределяем функцию   \varphi(x-ct)   при x < ct. Для этого подставим общее решение   u(x,t)=\varphi(x-ct)+\psi(x+ct)   в граничное условие u(0,t)=0. В результате получим   \varphi(-ct)=-\psi(ct)  ,откуда
\end{verbatim}

\[
\Delta u = 0
\]

Подстановка этого выражения для ф в общее решение дает

\[
\Delta u + \varepsilon u^2 = 0
\]

Комбинируя решения для x \textless{} ct и x \textgreater{} ct,
окончательно получаем (см. рис. 18.6)

\Дельта u + u\^{}2 = 0

\[
\Delta u = 0 \rightarrow \Delta u + \varepsilon u^2 = 0 \rightarrow \Delta u + u^2 = 0
\begin{cases} \frac{1}{2} \left[ f(x - ct) + f(x + ct) \right] + \frac{1}{2c} \int_{x - ct}^{x + ct} g(\xi) d\xi, & x \ge ct, \\ \frac{1}{2} \left[ f(x + ct) - f(ct - x) \right] + \frac{1}{2c} \int_{ct - x}^{x + ct} g(\xi) d\xi, & x < ct. \end{cases}
u_0 \rightarrow u = \sum_{k=0}^{\infty} \varepsilon^k u_k \rightarrow u = \sum_{k=0}^{\infty} u_k
\Delta u + cu^2 = 0,
\]

Этим завершается наша лекция. Интерпретация решения (18.8) будет
приведена в замечаниях.

\pandocbounded{\includegraphics[keepaspectratio]{images/_page_137_Figure_9.jpeg}}

\begin{verbatim}
Рис. 18.6. Интерпретация задачи Коши для полуограниченной струны в xt-плоскости a—возмущение отражается от границы;   \delta  —возмущение распространяется вдоль характеристики из точки (  x_0+ct_0  , 0).
\end{verbatim}

\section{ЗАУВАЖЕННЯ}\label{ux437ux430ux443ux432ux430ux436ux435ux43dux43dux44f-1}

\begin{verbatim}
1. Решение (18.8) полностью соогветствует нашим представлениям о полубесконечной струне с граничным условием u(0, t) = 0. При   x \geqslant ct   оно совпадает с формулой Даламбера для бесконечной струны. Если x < ct, то формула модифицируется таким образом, что в нее включается волна, отраженная от границы (при отражении знак волны меняется на противоположный).2. Если граничное условие u(0, t) = 0 изменить на другое, то и решение (18.8), естественно, изменится. Можно получить в явном виде решения для граничных условий
\end{verbatim}

\[
\Delta u + \varepsilon u^2 = 0
\]

или

\[
(46.1) u = u_0 + \varepsilon u_1 + \varepsilon^2 u_2 + \dots,
\]

Дополнительную информацию по этим вопросам читатель может найти в
литературе, указанной в конце книги.3. Прямые линии

\Дельта u + u\^{}2 = 0

\begin{verbatim}
x - ct = \text{const}  называются характеристиками. Возмущение распространяется вдоль характеристик. Понятие характеристики тесно связано с уравнениями гиперболического типа.
\end{verbatim}

\section{ЗАДАЧІ}\label{ux437ux430ux434ux430ux447ux456}

\begin{enumerate}
\def\labelenumi{\arabic{enumi}.}
\tightlist
\item
  Решите смещанную задачу для полубесконечной струны
\end{enumerate}

\[
\Delta u + u^2 = 0,
\Delta u = 0, \quad 0 < r < 1, \quad 0 \le \theta < 2\pi,
\]

Постройте графики решения для разл чных моментов времени.2. Щоб
розв'язати попередню задачу, можна використати наступний алгоритм:а)
продолжить начальные условия на всю действительную ось по формулам

\begin{verbatim}
имеет только вулевое решение. - Прим. ред.2) Здесь и ниже речь идет об отыскании одного из решений (второе таким методом не находится).— Прим ред.3) Эта «ракета», во-первых, не при всех граничных условиях долетает до є=1 (где при опредсленных вещественных граничных условиях открывается мир комплексных решений) и, во-вторых, ни при каких в не долетает до второго вещественного решения, когда оно существует.— Прим. ред.ции  u_0 ,  u_1 ,  u_2 , ..., входящие в степенной ряд (46.1). В оставшейся части лекции мы покажем, как этим методом можно решить различные краевые задачи.# Решение нелинейного уравнения \Delta u + u^2 = 0 методом возмущенийПредположим, что нам нужно решить следующую нелинейную задачу Дирихле:(46.2)
\end{verbatim}

\(u_t(x, 0) = 0, -\infty < x < 0.\) b) середнє дві хвилі, що рухаються
вліво і вправо (як це було в попередній лекції);\textbf{в}) рассмотреть
полученное решение при \(x \ge 0\) .Використовуючи цей алгоритм,
побудуйте розв'язок для різних помножених на t, якщо початкова умова
задана графічно.

\pandocbounded{\includegraphics[keepaspectratio]{images/_page_139_Figure_3.jpeg}}

\begin{enumerate}
\def\labelenumi{\arabic{enumi}.}
\setcounter{enumi}{2}
\tightlist
\item
  Розв'язати змішану задачу для напівнескінченної струни
\end{enumerate}

Розглянемо цю нелінійну задачу як збурену лінійну задачу:(46.3)

Розв'язок якої \(u_0(r, \theta) = r \cos \theta\) відомий. Як змінити
\(u_0(r, \theta)\) , щоб отримати рішення проблеми (46.2)? Як уже
згадувалося, розглянемо клас задач \(\Delta u + \epsilon u^2 = 0\) з
граничними умовами і шукаємо розв'язок кожної задачі з цього класу у
вигляді(46.4)

Воспользуйтесь методом, аналогичным рассмотренному в лекции. Дайте
физическую интерпретацию этой задачи.4. Предположим, что колебания
струны описываются уравнением \(u_{tt} = u_{xx}\),начальное смещение
изображено на рисунке, а начальная скорость равна нулю. Изобразите
решение этой задачи в xt-плоскости. Обратите внимание на то, что
начальные условия в этой задаче разрывны.

\pandocbounded{\includegraphics[keepaspectratio]{images/_page_139_Figure_8.jpeg}}

\section{ВОЛНОВОЕ УРАВНЕНИЕ И ГРАНИЧНЫЕ УСЛОВИЯ МЕТА ЛЕКЦІЇ: Познакомыть
читателя с основными типами граничных условий для волнового
уравнения}\label{ux432ux43eux43bux43dux43eux432ux43eux435-ux443ux440ux430ux432ux43dux435ux43dux438ux435-ux438-ux433ux440ux430ux43dux438ux447ux43dux44bux435-ux443ux441ux43bux43eux432ux438ux44f-ux43cux435ux442ux430-ux43bux435ux43aux446ux456ux457-ux43fux43eux437ux43dux430ux43aux43eux43cux44bux442ux44c-ux447ux438ux442ux430ux442ux435ux43bux44f-ux441-ux43eux441ux43dux43eux432ux43dux44bux43cux438-ux442ux438ux43fux430ux43cux438-ux433ux440ux430ux43dux438ux447ux43dux44bux445-ux443ux441ux43bux43eux432ux438ux439-ux434ux43bux44f-ux432ux43eux43bux43dux43eux432ux43eux433ux43e-ux443ux440ux430ux432ux43dux435ux43dux438ux44f}

Потім знаходимо розв'язок задачі (46.2), поклавши
\(\varepsilon = 1^{\circ}\) ).Вираз (46.4), який претендує на розв'язок,
буде замінений на обурену проблему(PPP)

на конечном отрезке.В дальнейшем будут рассматриваться три основных типа
граничных условий.1. Заданный режим (граничные условия 1-го рода)

, (GI) \(u(1, \theta) = \cos \theta\) .У результаті ми отримуємо

\(u(L, t)\) = \(g_2\)(t). 2. Заданные силы (граничные условия 2-го рода)

\[
u_0 (1, \theta) + \varepsilon u_1 (1, \theta) + \varepsilon^2 u_2 (1, \theta) + \ldots = \cos \theta.
\]

\(u_x\)(L, t) = \(g_2\)(t). 3. Упругое закрепление (граничные условия
3-го рода)

\[
P_{0}
\]

\begin{verbatim}
u_x(L, t) - \gamma_2 u(L, t) = g_2(t).  Могут встречаться и произвольные комбинации этих трех типоз граничных условий. В лекции будет показан физический смысл этих граничных условий.Итак, пока что мы рассмотрели единственный вид волнового движения—поперечные колебания струны. Читатель должен понимать, что это только один из видов волновых движений.Ниже приведен перечень основных типов волн:- 1. Звуковые волны (продольные волны).
\end{verbatim}

\begin{itemize}
\tightlist
\item
  \begin{enumerate}
  \def\labelenumi{\arabic{enumi}.}
  \setcounter{enumi}{1}
  \tightlist
  \item
    Электромагнитные волны, в том числе световые.
  \end{enumerate}
\item
  \begin{enumerate}
  \def\labelenumi{\arabic{enumi}.}
  \setcounter{enumi}{2}
  \tightlist
  \item
    Колебания твердых тел (продольные, поперечные и крутильные).
  \end{enumerate}
\item
  \begin{enumerate}
  \def\labelenumi{\arabic{enumi}.}
  \setcounter{enumi}{3}
  \tightlist
  \item
    Волны вероятности в квантовой механике.
  \end{enumerate}
\item
  \begin{enumerate}
  \def\labelenumi{\arabic{enumi}.}
  \setcounter{enumi}{4}
  \tightlist
  \item
    Волны на воде (поперечные волны).
  \end{enumerate}
\item
  \begin{enumerate}
  \def\labelenumi{\arabic{enumi}.}
  \setcounter{enumi}{5}
  \tightlist
  \item
    Колебания струны (поперечные волны).Цель нашей лекции --- обсудить
    различные типы граничных условий, которые возникают при решении
    физических задач волнового характера. Мы остановимся только на
    одномерных задачах и линейных граничных условиях. Обычно различают
    граничные условия трех типов1. Заданные режимы в граничных точках
    (граничные условия 1-го рода)
  \end{enumerate}
\end{itemize}

\[
\qquad u_{\theta}(r, \theta) = r \cos \theta,
\]

\(u(L, t)\) = \(g_2\)(t). 2. Заданные силы в граничных точках (граничные
условия 2-го рода)

\[
\begin{cases} \Delta u_{1} = -u_{0}^{2}, \\ u_{1}(1, \theta) = 0, \\ u_{2}(1, \theta). \end{cases}
\]

\(u_x\)(L, t) = \(g_2\)(t). 3. Упругое закрепление в граничных точках
(граничные условия 3-го рода)

\[
\begin{cases} \Delta u_{2} = -2u_{0}u_{1}, \\ u_{2}(1, \theta). \end{cases}
\]

\begin{verbatim}
u_x(L, t) - \gamma_2 u(L, t) = g_2(t).  Мы начнем с граничных условий 1-го рода.1. Заданные режимы в граничных точках.Нам предстоит решить задачу вида(19.1)
\end{verbatim}

\[
P_{1}
\]

Щоб розв'язати \emph{неоднорідне рівняння}, ми також звернемося до
методу, який використовувався в теорії звичайних диференціальних рівнянь
і складається з наступного:- 1) знайти загальний розв'язок \(y_n\)
однорідного рівняння;

\begin{itemize}
\tightlist
\item
  \begin{enumerate}
  \def\labelenumi{\arabic{enumi})}
  \setcounter{enumi}{1}
  \tightlist
  \item
    знайти часткове розв'язок \(y_p\) гетерогенного рівняння;
  \end{enumerate}
\end{itemize}

\begin{verbatim}
- 3) Підставити y_h + y_p у початкову умову та визначити константи.Ми використаємо цей метод для розв'язання нашої проблеми. У нашій задачі P_1 загальний розв'язок однорідного рівняння \Delta u=0 має вигляд із розділеними змінними
\end{verbatim}

\[
\begin{aligned} u_{tt} &= c^2 u_{xx}, & 0 < x < 1, & 0 < t < \infty, \\ u(0, t) &= g_1(t), & 0 < t < \infty, \\ u(1, t) &= g_2(t), & 0 < t < \infty, \\ (HY) &\begin{cases} u(x, 0) &= f(x), \\ u_t(x, 0) &= g(x), \end{cases} & 0 \leqslant x \leqslant 1, \end{aligned}
\]

Часткове (лише одне) розв'язок неоднорідного рівняння

Ми розглядаємо у формі

\[
u_x(0, t) = 0, \quad u_x(L, t) = 0.
\]

Правильні частини форми \(r^n\) , \(r^n \cos{(n\theta)}\) ,
\(r^n \sin{(n\theta)}\) ведуть до рішень \(Ar^{n+2}\) ,
\(Br^{n+2} \cos{(n\theta)}\) \(Cr^{n+2} \sin{(n\theta)}\)
відповідно.Підставляючи функції \(u_p(r, \theta)\) у гетерогенне
рівняння, отримуємо A = -1/32, B = -1/24.Отже,

\[
u_{x}(1, t) = \frac{1}{k} v(t) - (k - \text{модуль Юнга}).
\]

Щоб розв'язати задачу, \(P_1\) має зробити останній крок --- підставити
загальне розв'язання \(u(r, \theta) = u_h(r, \theta) + u_p(r, \theta)\)
на граничні умови \(u(1, \theta) = 0\) і знайти коефіцієнти \(a_n\) і
\(b_n\) .Після цього ми отримуємо

\[
u_{x}(0, t) = \frac{h}{T} u(0, t),
\]

Отже, \(a_0=1/32\) , \(a_2=1/24\) , і всі інші \(a_n\) і \(b_n\)
дорівнюють нулю. Отже, рішення проблеми \(P_1\) знайдено:

\[
u_{x}(L, t) = -\frac{h}{T} u(L, t).
\]

Функція \(u_1(r, \theta)\) називається пертурбацією першого порядку до
\(u_4(r, \theta)\) . Додавши \(u_0\) і \(u_1\) , отримуємо новий підхід
до розв'язання задачі (46.2):(46.5)

\[
u_{x}(0, t) - \frac{h}{T}u(0, t) = 0,
\]

Щоб знайти збурення наступного порядку \(u_{z}(r, \theta)\) , потрібно
розв'язати задачу

\[
u_{x}(L, t) + \frac{h}{T}u(L, t) = 0.
\begin{cases} \Delta u_{1} = -2u_{0}u_{1}, \\ u_{2}(1, \theta) = 0, \end{cases}
\]

де \(u_{\phi}(r,\theta)\) і \(u_{1}(r,0)\) вже відомі, раніше знайдені
функції. Звісно, виконувати всі ці алгебраїчні перетворення без допомоги
комп'ютера надзвичайно складно. На щастя, вираз (46.5) є майже точним
рішенням нашої проблеми. Насправді, якщо підставити цей вираз зліва від
\(\Delta u + u^{2}\) нашого нелінійного рівняння, отримаємо вираз, який
майже всюди відрізняється від нуля в межах кола 0 \textless{} r
\textless{} 1.Окрім розв'язання нелінійних задач, теорія збурень може
використовуватися для розв'язання задач у областях з неправильними
межами (звісно, якщо нерегулярність не дуже велика). Давайте розглянемо
простий приклад.\# Приклад задачі з порушеною межеюОсквернення можна
знищити, не позбавляючи його права. Можна знайти розвиток рівня Лаплас
посередині деформованого стовпа, оскільки пробурив рівень Лаплас у
центрі круглої області. У застосуванні дозволено знайти потенціал у
квадраті за кривою \(r=1+\frac{1}{4}\sin\theta\) (за полярними
координатами). Нехай потенціал у на межі заданий. Тобто, нам потрібно
вирішити проблему(46.6)

\[
u \left( 1 + \frac{1}{4} \sin \theta, \ \theta \right) = \cos \theta, \quad 0 \le \theta \le 2\pi.
\begin{aligned} u_x(0, T) &= \frac{h}{T} [u(0, t) - \theta_1(t)], \\ u_x(L, T) &= -\frac{h}{T} [u(L, t) - \theta_2(t)]. \end{aligned}
u\left(1+\frac{1}{4}\sin\theta,\ \theta\right)=\cos\theta
u(1+\epsilon\sin\theta, \theta) = \cos\theta, \quad 0 \le \epsilon \le 1/4.
\]

Этим мы завершаем обсуждение наиболее важных граничных условий,
связанных с гиперболическими задачами. Следующие несколько лекций мы
посвятим решению задач, содержащих полобные граничные условия.

\section{ЗАУВАЖЕННЯ}\label{ux437ux430ux443ux432ux430ux436ux435ux43dux43dux44f-2}

\begin{enumerate}
\def\labelenumi{\arabic{enumi}.}
\tightlist
\item
  В лекции не рассматривалось граничное условие еще одного типа, когда
  струна подвергается действию силы, пропорциональной скорости и
  направленной в противоположном направ-лении. Граничное условие такого
  типа (на левом конце) записывается в виде
\end{enumerate}

\[
f(x+h) = f(x) + f'(x)h + f''(x)\frac{h^2}{2!} + \dots,
\]

\begin{enumerate}
\def\labelenumi{\arabic{enumi}.}
\setcounter{enumi}{1}
\tightlist
\item
  Нелинейное упругое закрепление левого конца приводит к граничному
  условию вида
\end{enumerate}

\[
u(1+\varepsilon\sin\theta,\ \theta)=u(1,\ \theta)+u_r(1,\ \theta)(\varepsilon\sin\theta)+u_{rr}(1,\ \theta)\frac{(\varepsilon\sin\theta)^2}{21}+\ldots
\]

\begin{verbatim}
где   \varphi(u)   — произвольная функция величины u. Например, граничное условие
\end{verbatim}

\[
\Delta u = 0, \quad 0 < r < 1 + \varepsilon \sin \theta,
\]

говорит о том, что возвращающая сила на левом конце пропорциональна кубу
смещения (а не первой степени и, как это будет согласно закону Гука в
линейном случае).3. Если к нижнему концу пружины, совершающей продольные
колебания, прикрепить массу \emph{m}, то граничное условие следует
задавать в виде

Коли всі розрахунки будуть виконані, ми внесемо \(\varepsilon=1/4\) у
кінцеві результати. Звісно, здається, що цю задачу не можна легко
розв'язати, але її можна звести до послідовності задач, у кожній з яких
визначена одна з функцій \(u_{\mathrm{ }e}\) , \(u_{\mathrm{ }e}\) ,
\ldots, що включені у розклад

\section{ЗАДАЧІ Опираясь на интуитивные представления о граничных
условиях различных типов, изобразить в общих чертах решение
задачи}\label{ux437ux430ux434ux430ux447ux456-ux43eux43fux438ux440ux430ux44fux441ux44c-ux43dux430-ux438ux43dux442ux443ux438ux442ux438ux432ux43dux44bux435-ux43fux440ux435ux434ux441ux442ux430ux432ux43bux435ux43dux438ux44f-ux43e-ux433ux440ux430ux43dux438ux447ux43dux44bux445-ux443ux441ux43bux43eux432ux438ux44fux445-ux440ux430ux437ux43bux438ux447ux43dux44bux445-ux442ux438ux43fux43eux432-ux438ux437ux43eux431ux440ux430ux437ux438ux442ux44c-ux432-ux43eux431ux449ux438ux445-ux447ux435ux440ux442ux430ux445-ux440ux435ux448ux435ux43dux438ux435-ux437ux430ux434ux430ux447ux438}

Після підстави ряду в (467) отримуємо таку послідовність задач для
визначення функцій \(u_a\) , \(u_1\) , \ldots;

\[
\begin{cases} \Delta u_{0} = 0, & 0 < r < 1 \text{ (внутри круга),} \\ u_{0}(1, \theta) = \cos \theta, & u_{0}(r, \theta) = r \cos \theta, \end{cases}
\]

\begin{verbatim}
\begin{array}{ll} (\mathrm{Y} \, \Pi) & u_{tt} = c^2 u_{xx}, \quad 0 < x < 1, \quad 0 < t < \infty, \\ (\mathrm{\Gamma} \, \mathrm{Y}) & \left\{ \begin{array}{ll} u \, (0, \, t) = 0, \\ u \, (1, \, t) = \sin t, \end{array}
\end{verbatim}

\[
\begin{cases} \Delta u_{1} = 0, & 0 < r < 1 \text{ (внутри круга),} \\ u_{1}(1, \theta) = -\sin \theta \frac{\partial u_{0}(1, \theta)}{\partial r} = -\sin \theta \cos \theta \end{cases}
u = u_0 + \frac{1}{4}u_1 + \left(\frac{1}{4}\right)^2 u_1 + \cdots
u_u + \frac{1}{4} u_t
\begin{array}{ll} u \, (x, \, 0) = 0, \\ u_{t} \, (x, \, 0) = 0, \end{array}
u_t = (1+x) u_{xx}, \\ u(x, 0) = q(x), \quad -\infty < x < \infty,
\]

\begin{verbatim}
\right. & 0 \leqslant x \leqslant 1, \end{array}
\end{verbatim}

Якщо розглянути рівняння з параметром(46.11)

\[
u_t = (1 + \varepsilon x) u_{xx},
\]

і шукати його розв'язок у вигляді

\[
u = u_0 + \varepsilon u_1 + \varepsilon^2 u_2 + \cdots,
\]

то після підстановки отримуємо послідовність задач

\[
P_0:\ \begin{cases}
\dfrac{\partial u_0}{\partial t} = \dfrac{\partial^2 u_0}{\partial x^2}, \\
u_0(x,0) = \varphi(x),
\end{cases}
\qquad P_1:\ \begin{cases}
\dfrac{\partial u_1}{\partial t} - \dfrac{\partial^2 u_1}{\partial x^2} = x\,\dfrac{\partial^2 u_0}{\partial x^2}, \\
u_1(x,0) = 0.
\end{cases}
\]

\begin{verbatim}
Зауважимо, що в перших двох задачах коефіцієнти сталі. Параметр \varepsilon має бути достатньо малим, інакше нескінченний ряд може виявитися розбіжним.#### ЗАДАЧИ

1. Подставьте разложение (46.8) в задачу (46.7) и получите последовательность задач $P_0,\ P_1,\ P_2,\ \dots$.

2. Покажите, что нелинейную задачу

$$
\begin{aligned}
\Delta u + u^2 &= 0, && 0 \leqslant r < 1,\ 0 \leqslant \theta < 2\pi, \\
u(1,\theta) &= \cos\theta, && 0 \leqslant \theta < 2\pi,
\end{aligned}
$$

можно свести к последовательности линейных задач $P_0,\ P_1,\ P_2$.

3. Подставьте приближенное решение (46.5) в нелинейную задачу 2 и оцените его точность.

4. Решите задачу $P_1$ из примера с возмущенной границей и проверьте, насколько хорошо приближенное решение

$$
u(r,\theta) = u_0(r,\theta) + \frac{1}{4}u_1(r,\theta)
$$

удовлетворяет соотношениям

$$
\Delta u = 0,
u\left(1 + \frac{1}{4}\sin\theta,\ \theta\right) = \cos\theta.
$$

# Лекція 47# РОЗВ'ЯЗАННЯ ДИФЕРЕНЦІАЛЬНИХ РІВНЯНЬ У ПОХІДНИХ ПОХІДНИХ ЗА ДОПОМОГОЮ МЕТОДУ КОНФОРМНИХ ВІДОБРАЖЕНЬМЕТА ЛЕКЦІЇ: Показати, як деякі двовимірні граничні задачі можна перетворити на інші, простіші задачі за допомогою конформного відображення. Припустимо, наприклад, що ми хочемо розв'язати рівняння Лапласа 

$$
\begin{cases} u(0, t) = 0, \\ u_x(1, t) = -u(1, t), \\ u(x, 0) = x, \\ u_t(x, 0) = 0. \end{cases}
$$

 з певною граничною умовою в області комплексної форми (площина змінних x і y). Цю задачу можна перетворити на нову, в якій потрібно знайти розв'язок рівняння Лапласа 

у простішій області (площина змінних і і і).

![](images/_page_147_Picture_5.jpeg)

г. 20.1. Три типові стоячі хвилі: $X(x)$ $T(t)$.Декілька прикладів стоячих хвиль показано на рис. 20.1. Якщо профілі цих стоячих хвиль відомі  $X_n$(x)  і природа коливань кожної хвилі  $T_n$(t)  , то розв'язок проблеми коливань гітарних струн є суперпозицією найпростіших коливань  $X_n$(x)\,$T_n$(t)  ісывается в виде

$$
= \ln\left[\frac{z-1}{z+1}\right] + i \arg\left[\frac{z-1}{z+1}\right]
$$
\end{verbatim}

Залишилося лише обрати шанси c\_n так, щоб при t = 0 початкові умови
u\left(x,~0\right)=f\left(x\right) і
u\_t\left(x,~0\right)=g\left(x\right) були виконані.Задачу о колебаниях
гитарной струны мы решим методом разделения переменных.\# Розв'язання
проблеми коливань обмеженої струни методом розділення зміннихЧтобы
решить смешанную задачу(УЧП) \textsubscript{}\textasciitilde{}

\[
v = \arg\left[\frac{z-1}{z+1}\right] = \arg\left[\frac{x+iy-1}{x+1+iy}\right] =
= \arg\left[\frac{2y}{x^2+y^2-1}\right]
\varphi(x, y) = \frac{1}{\pi} \arctan\left[\frac{2y}{x^2 + y^2 - 1}\right].
\begin{cases} u(0, t) = 0, \\ u(L, t) = 0, \end{cases}
x^{2} + y^{2} = 1,
\]

якщо потенціал на внутрішньому колі дорівнює одинці, а на зовнішньому
колі дорівнює двом (рис. 47.5). Іншими словами, нам потрібно розв'язати
проблему(47.3)

Проблема полягає в тому, щоб знайти конформне відображення, яке
перетворює цю область на щось простіше, де нашу проблему легко
розв'язати. У цьому випадку очевидно, що ця проста ділянка має бути
кільцем. Відповідна експозиція Р И С У Н О

\pandocbounded{\includegraphics[keepaspectratio]{images/_page_149_Figure_5.jpeg}}

20.2. Построение стоячих воли при раздичных значениях параметра
λ.водными и граничным условиям. Затем останется только найти такую
линейную комбинацию этих колебаний, чтобы сумма удовлетворяла начальным
условиям при t=0.Подстановка (20.2) в граничные условия
\(u(0, t) = u(L, t) = 0\) дает

\[
u(0, t) = X(0)T(t) = D\bigl[A\sin(\alpha\beta t) + B\cos(\alpha\beta t)\bigr] = 0 \Rightarrow D = 0,
\]

и

\[
u(L, t) = X(L)T(t) = C\sin(\beta L)\bigl[A\sin(\alpha\beta t) + B\cos(\alpha\beta t)\bigr] = 0 \Rightarrow \sin(\beta L) = 0.
\]

Другими словами, константа разделения \(\beta\) должна удовлетворять
уравнению \(\sin(\beta L) = 0\), откуда

\[
\beta_n = \frac{n\pi}{L}, \qquad n = 0, 1, 2, \ldots
\]

Заметим, что если во втором уравнении (20.3) положить \(C=0\), то
получится тривиальное решение \(X(x)T(t) \equiv 0\). Следовательно, мы
нашли последовательность элементарных колебаний струны:

\[
u_n(x, t) = X_n(x)T_n(t) = \sin\!\left(\frac{n\pi x}{L}\right)\left[a_n \sin\!\left(\frac{n\pi \alpha t}{L}\right) + b_n \cos\!\left(\frac{n\pi \alpha t}{L}\right)\right].
\]

(20.4)або

\[
u_n(x, t) = R_n \sin\!\left(\frac{n\pi x}{L}\right)\cos\!\left(\frac{n\pi \alpha (t - \delta_n)}{L}\right).
\]

Тут \(a_n\), \(b_n\), \(R_n\) і \(\delta_n\) --- довільні сталі. Кожне
таке елементарне коливання є стоячою хвилею.

\pandocbounded{\includegraphics[keepaspectratio]{images/_page_114_Figure_4.jpeg}}

Рис. 15.3 Конвективная волна a---передиції фронт примеси, \(\delta\)
---на одну единицу вперед волнового фронтаОбратим внимание на то, что в
задаче (15.2) мы отодвинули границу в --- \(\infty\) (теперь решается
задача Коши), так что она не оказывает теперь влияние на оценку решение.
Чтобы решить задачу (15.2), можно воспользоваться преобразованием
Лапласа Р и с

\pandocbounded{\includegraphics[keepaspectratio]{images/_page_114_Figure_7.jpeg}}

. 15 4 Начальные условия в задаче конвективной диффузии.ло переменной t
или преобразованием Фурье по переменной x. Однако в этом случае гораздо
интереснее поступить совсем по-другому Введем новую систему координат,
которая движется вдоль старой со скоростью V. Другими словами, вместо
системы координат, привязанной к берегу реки, мы рассмотрим систему
координат, которая движется со скоростью фронта примеси (конечно, при
налични диффузии фронт будет «смазан»). С точки зрения математики это
означает, что мы заменяем пространственную координату x на новую
\(\xi = x - Vt\) .Теперь ясно, чтокогда \(\xi = 0\) , мы находимся па
фронте распространяющейся примеси.когда \(\xi = 1\) , мы находимся на
одну единицу длины впереди фронта.когда \(\xi = -1\) , мы находимся на
одну единицу длины позади фронта.Наша задача --- преобразовать исходную
задачу с НУ

\[
\overline{D} = A\xi_{xx} + B\xi_{xy} + C\xi_{yy} + D\xi_{x} + E\xi_{y},
[\xi_x/\xi_y] = -B/2A
\xi(x, y) = c
\frac{dy}{dx} = -\left[\xi_x/\xi_y\right] = B_t 2A.
u_{\alpha\alpha} - u_{\beta\beta} = \frac{-\beta u_{\alpha} - \alpha u_{\beta}}{2\alpha\beta}.
\overline{B} = 2A\xi_x\eta_x + B(\xi_x\eta_y + \xi_y\eta_x) + 2C\xi_y\eta_y,
\alpha = \xi + \eta = (y^2 - x^2) + (y^2 + x^2) = 2y^3,
\]

1) Завдяки рівностіІ з \(B^2 - 4AC = 0\) році , то

\[
\beta = \xi - \eta = (y^2 - x^2) - (y^2 + x^2) = -2x^2.
\]

Оскільки
\(\xi_x/\xi_y = -B \ 2A = -2 \sqrt{A\tilde{C}}/2A = -\sqrt{C/A}\)
останній вираз для B буде записаний як

\[
u_{\xi\xi} - u_{\eta\eta} = \Psi(\xi, \, \eta, \, u_{\xi}, \, u, u_{\eta}).
\]

І оскільки \(\overline{A}\) дорівнює нулю, то \(\overline{B}\) , звісно,
також дорівнює нулю.Оскільки при нашому виборі \(\xi\) обидва
коефіцієнти \(\overline{A}\) і \(\overline{B}\) повертаються до нуля, ми
можемо взяти будь-яку функцію як \(\eta\) (якщо вона не пропорційна
координаті \(\xi\) ). У прикладі можна вибрати \(\eta = y\) .Залишається
знайти канонічну форму рівняння в нових координатах. Для цього потрібно
просто підставити формули \(\xi\) і \(\eta\) (41.3) і знайти всі шанси
\(\overline{A}\) , \(\overline{B}\) , \(\overline{C}\) ,
\(\overline{D}\) , \(\overline{E}\) , \(\overline{F}\) і
\(\overline{G}\) . На цьому вивчення параболічного випадку завершується,
але перш ніж перейти до іншої теми, розглянемо простий приклад.\#
Перетворення параболічного рівняння \(u_{xx} + 2u_{xy} + u_{yy} = 0\) у
канонічну формуРозглянемо рівняння

\[
3u_{xx} + 7u_{xy} + 2u_{yy} = 0
\]

Його коефіцієнти мають такі значення: A=1, B=2, C=1, D=E=F=G=0. Отже,
\(B^2-4AC=0\) для всіх значень x і y. Щоб знайти нові координати \(\xi\)
і \(\eta\) та канонічну форму рівняння, ми рухаємося наступним
чином.КРОК 1. Запишемо \emph{характеристичне рівняння} (тепер лише одне,
а не два)

\[
u_{\xi\eta} = \Phi(\xi, \eta, u, u_1, u_2).
\]

Розв'язуючи це рівняння відносно y, отримуємо

\[
u_{\alpha\alpha} - u_{\beta\beta} = \Psi(\alpha, \beta, u, u_{\alpha}, u_{\beta}).
\]

,Де --- характеристика \(\xi = y - x\) визначена. Використання цієї
характеристики забезпечує досягнення рівності \(\overline{A} = 0\) .
Координату \(\eta\) можна вибрати довільно (якщо вона не залежить від
\(\xi\) ). Ми оберемо його наступним чином:

\[
u_{xx} + 4u_{xy} = 0.
\]

.Нові координати

\[
u_{tt} = e^{2} (u_{xx} + u_{yy} + u_{zz}),
\]

изображены на рис. 41 1.

\pandocbounded{\includegraphics[keepaspectratio]{images/_page_117_Figure_3.jpeg}}

Рис. 15.6. Диффузня из области высокой концентрации в область низкой
концентрации. Чем больше коэффициент диффузии, тем быстрее установится
стационарное значение.Другими словами, в зависимости от относительной
величины D (коэффициента диффузии) и V (скорости потока) решение
движется слева направо со скоростью V и в то же самое время передний
фронт расплывается со скоростью, определяемой величиной D.

\pandocbounded{\includegraphics[keepaspectratio]{images/_page_117_Figure_6.jpeg}}

Рис. 15.7. Решение задачи конвективнои диффузии. Вещество одновременно
движется и диффундирует.(На рис. 15.7 показано, как расплывается
передний фронт концентрации.) \# ЗАУВАЖЕННЯ Преобразование координат
является важным методом решения уравнений с частными производными.
Выбрав подходящую систему координат, можно существенно упростить
уравнение. \# ЗАДАЧІ 1. Решите задачу Коши:

Але тепер у B\^{}2-4AC\textless0 . Переходячи до нових пояснювальних
змінних, ми хочемо перетворити її у нову форму

\begin{enumerate}
\def\labelenumi{\arabic{enumi}.}
\setcounter{enumi}{1}
\tightlist
\item
  Найдите решение задачи Коши для уравнения конвективной диффузии:
\end{enumerate}

\begin{verbatim}
Це можливо, але лише шляхом введення складних координат. Щоб знайти ці комплексні координати  \xi  і  \eta  , розв'яжемо характеристичні рівняння, так само як у гіперболічному випадку:
\end{verbatim}

Воспользуйтесь преобразованием из лекции 8.3. Найдите решение задачи
переноса:

\[
u_{tt} = c^2 (u_{xx} + u_{yy}), \quad
\frac{dy}{dx} = \sqrt{-4x^2} = 2ix
\\ -\infty < x < \infty, \quad 0 < t < \infty,
u_t = Du_{xx} - Vu_x, \quad -\infty < x < \infty, \quad 0 < t < \infty,
u(x, 0) = e^{-x^2}, \quad -\infty < x < \infty.
y^2 u_{xx} + x^2 u_{yy} = 0,
\frac{dy}{dx} = \frac{B - \sqrt{B^2 - 4AC}}{2A} = -\frac{\sqrt{-4x^2y^2}}{2y^2} = -i\frac{x}{y},
\begin{cases} -\infty < x < \infty, \\ -\infty < y < \infty, \\ -\infty < z < \infty, \end{cases}
u(\xi, \tau) = \frac{1}{2\sqrt{D\pi\tau}} \int_{-\infty}^{+\infty} e^{-\beta^2} e^{-(\xi-\beta)^2/4D\tau} d\beta.
u_{tt} = \alpha^2 u_{xx}
P_{2}
u_{\xi\eta} = \psi(\xi, \eta, u, u_{\xi}, u_{\eta}).
\begin{aligned} u_{tt} &= c^2 \Delta u, \\ (HY) & \begin{cases} u(x, y, z, 0) = 0, \\ u_t(x, y, z, 0) = \psi(x, y, z), \end{cases} \end{aligned}
\xi(x, y) = y^2, \quad \eta(x, y) = x^2.
T \sin \theta_2 - T \sin \theta_1 =
P_{1}
\Delta x \rho u_{tt} = T \left[ u_x \left( x + \Delta x, \ t \right) - u_x \left( x, \ t \right) \right] + \Delta x F \left( x, \ t \right) - \frac{1}{\Delta x \rho u_t} \left( x, \ t \right) - \Delta x \gamma \dot{u} \left( x, \ t \right)
\]

где \(\overline{\psi}\) --- \emph{среднее} значение начального
распределения \(\psi\) \emph{по сфере} радиуса ct с центром в точке (x,
y, z), т. е.

\[
u_{tt} = \alpha^{3} u_{xx} - \beta u_{t} - \gamma u + F(x, t).
\]

\pandocbounded{\includegraphics[keepaspectratio]{images/_page_134_Figure_7.jpeg}}

Рис. 18.4. Решение задачи (18.4) в xt-плоскости.Задача (18.4) описывает
поведение струны, которой сообщается рачальная единичная скорость на
отрезке -1 \textless{} x \textless{} 1. Смещение . .описывается формулой
Даламбера

\[
\begin{aligned} & (\text{УЧП}) & \phi_{xx} + \phi_{yy} = 0, & \text{внутри } D, \\ & (\text{ГУ}) & \begin{cases} \varphi(x, y) = 1 & \text{на } x^2 + y^2 = 1, \\ \varphi(x, y) = 2 & \text{на } (x - 1)^2 + y^2 = 9. \end{cases} \end{aligned}
\]

\begin{verbatim}
является теперь функцией коэффициентов   a_1a_2\dots a_n  . Следовательно, для того чтобы найти минимум функционала J, приравняем нулю частные производные
\end{verbatim}

\[
Aa = b
b_i = -\int_0^1 \int_0^1 f(x, y) \varphi_i(x, y) dx dy,
\]

где

Это решение в различные моменты времени изображено на рис. 18.5.

\pandocbounded{\includegraphics[keepaspectratio]{images/_page_135_Figure_5.jpeg}}

\begin{verbatim}
Рис. 18.5. Решение задачи (18.4) в различные моменты времени.Этим завершается наша интерпретация формулы Даламбера в плоскости переменных   \boldsymbol{x}   и t. В оставшейся части лекции мы займемся решением смешанной задачи для полубесконечной струны:(18.6)
\end{verbatim}

\[
u = \gamma [(x-s)(x-t)-\gamma^2],
\]

\begin{verbatim}
Единственное отличие в том, что аппроксимирующие функции   \phi_1, \phi_2, \phi_3, \ldots, \phi_n   должны удовлетворять граничным условиям:
\end{verbatim}

\[
\gamma = 2t/[(x-t)^2 + y^2].
\]

\begin{verbatim}
является иравнение Лапласа, но это доказательство проводится аналогично тому, как это делалось в лекции 44 для функционала, зависящего от функции одной переменной.3. В книгах по вариационному исчислению показано, как строить энергетический функционал не только для уравнения Пуассона с граничным условием Дирихле u=0, но и для многих других дифференциальных уравнений и многих других типов граничных условий. Следовательно, большое число краевых задач можно решить, минимизируя соответствующие энергетические функционалы.4. Чем большее значение п мы выбираем, тем меньшее значение   J[\bar{u}_n]   получается. Значит, с увеличением n растет точность л приближенного решения дифференциального уравнения. Один из методов определения и, основан на последовательном вычислении функционала   J[u_n]   для возрастающих значений n. Вычисления прекращаются, если с ростом п функционал перестает практически уменьшаться.5. Если и велико, то проведение вычислений по методу Ритца требует привлечения ЭВМ. На рис. 45.2 (с. 345) приведена блок-схема решения краевой задачи методом Ритца.Пользователь должен позаботиться о подпрограмме, которая должна вычислять значения пробных функций. Такая подпрограмма может иметь вид:SUBROUTINE BC (X.Y.PHI)- SUBROUTINE PROVIDED BY THE USER TO C - EVALUATE THE - FUNCTIONS PHI(1), PHI(2), ..., PHI(N) DIMENSION PHI(20)PHI(1) = X\*Y\*(1-X)\*(1-Y) PHI(2) = X^*PHI(1)   (используемые функции)PHI(3) = Y\*PHI(1)PHI(N) = (что получится)RETURN END# **ЗА**ДАЧИНайдите энергетический функционал для задачи(УЧП)
\end{verbatim}

Если теперь преобразовать начальные условия краевой задачи, то получим
начальные условия для обыкновенного дифференциального уравнения

, 0 \textless{} x \textless{} 1 , 0 \textless{} y \textless{} 1 , (ГУ) u
= 0 на границе.2. Как минимизировать функционал

методом Ритца?УКАЗАНИЕ. Если ввести новую функцию \(z(x)\):

то можно заметить, что она удовлетворяет граничным условиям \(z(0)\) = 0
и \(z(1)\) = 0.1- 3. Напишите программу для проведения вычислений по
блоксхеме, изображенной на рис. 45 2.

\begin{itemize}
\tightlist
\item
  \begin{enumerate}
  \def\labelenumi{\arabic{enumi}.}
  \setcounter{enumi}{3}
  \tightlist
  \item
    Покажите, что уравнение
  \end{enumerate}
\end{itemize}

\[
\varphi(r) = a \ln r + b, \quad r = u^2 + v^2.
\]

является уравнением Эйлера --- Лагранжа для функционала

\[
\varphi(x, y) = 0.57 \ln(u^2 + v^2) + 1,
\]

\begin{enumerate}
\def\labelenumi{\arabic{enumi}.}
\setcounter{enumi}{4}
\tightlist
\item
  Задачу Дирихле
\end{enumerate}

Итак, решим теперь новое семейство задач Кощи:(ОДУ)

\begin{verbatim}
где A=0.06 и B=0.04. Как найти потенциальную энергию этого решения? Мы рекомендуем читателю вспомнить синуспреобразование и самостоятельно получить решение   u\left(x,y\right)  .<sup>\*)</sup> При такой замене возникает особенность в т. x=1. Лучше ввести функцию z(x) по формуле z(x)=y(x)=x. Прим. ред.
\end{verbatim}

\pandocbounded{\includegraphics[keepaspectratio]{images/_page_144_Picture_4.jpeg}}

Рис. 19.6. Возникновение неоднородных граничных условий упругого
закрепления: h---коэффициент упругости пружины.отрицательна. Однородные
граничные условия (19.2) можно записать в виде(19.3)

\[
P_{2}
\]

\(f(x)\) * \(g(x)\)

\[
\Delta u = 0, \quad 0 < r < 1 + \frac{1}{4} \sin \theta,
(1+x^2)^{-1}
\]

досягається в кінцевій точці сімейства граничних умов виду (див. рисунок
46 3)

\[
Tu_x(0, t) = -\beta u_t(0, t).
Tu_{x}(0, t) = \varphi[u(0, t)],
\frac{a}{x^2+a^2}
Tu_{x}(0, t) = -hu^{3}(0, t)
\frac{2ax}{(x^2+a^2)^2}
\]

mu\_\{tt\}(L, t) = -ku\_x(L, t) + mg.\$\$ \(P_{0}\)

\end{document}
